% ============================================================
% PAPER 3 of 11 -- computational_certification
% Assembled 2026-09-29 by content-and-merits partition.
% Title: Computational Certification of Obstruction
% Source: arena agent 1/paper rewrites/latex/paper2_computational_certification_v20.tex
% Partition note: Single-source. The computational question: adjoint safety rows, the continuous-to-finite theorem, complexity.
%% STATUS: structurally assembled as a NEW version. Editorial integration
%%         (de-duplication of shared framework, sibling cross-citations)
%%         is NOT yet done. See PAPERS_MANIFEST.md.
%% ============================================================

%% v12 (2026-09-29): added section \ref{priorart} and ten bibliography entries, all
%%   VERIFIED against publisher records before use.
%%   HONEST BASELINE: the two-sided SANDWICH IS NOT NEW -- Maidens et al. 2013 already
%%   produce a guaranteed under-approximation plus a free over-approximation used to
%%   bound the error. That is conceded in the text rather than claimed. The claimed
%%   difference is the object (a safety VALUE over an information state, against every
%%   measurable policy, rather than a kernel as a set of states) and the approximations
%%   (moment relaxation of information-adapted controls + adversarial scenarios, rather
%%   than a gridded state space).
%%   CITATION TRAP AVOIDED: the Automatica 2013 paper is MAIDENS, KAYNAMA, MITCHELL,
%%   OISHI and DUMONT -- Mitchell is THIRD author, not first. Do not 'correct' it.
%%   NOTATION COLLISION FLAGGED: beta in the error term delta_beta is deterministic,
%%   while beta in the scenario literature is a confidence parameter. Stated, not
%%   silently resolved.
%% PRIOR ART REQUIRED -- bar item, inserted 2026-09-29.
%% No related-work position, and this paper carries the fewest formal results of
%% the theory cluster (5). It must engage:
%%   - barrier certificates: Prajna and Jadbabaie (2004); Prajna, Jadbabaie and
%%     Pappas (2007); Prajna and Rantzer (2005); Maghenem and Sanfelice (2019).
%%   - Saint-Pierre (1994): kernel approximation and its complexity.
%%   - Hamilton-Jacobi reachability: Mitchell, Bayen and Tomlin (2005); Margellos
%%     and Lygeros (2011). Grid solvers are confined to five or six state
%%     dimensions, which is the natural computational comparator for any
%%     complexity claim made here.
%%   - Aubin and Catte (2002): bilateral fixed points for kernels and basins.
%% Claim to establish: the continuous-to-finite theorem and the complexity
%% statements either improve on or are provably incomparable with each of these.
%% Name the nearest neighbour.

% SIBLING PAPERS: Paper 1 (Obstruction Certificates under Incomplete Observation), Paper 2 (Exact Probabilistic Sufficiency), Paper 4 (Minimax Dual Certificates), Paper 5 (Exact Belief Computation).
% Cross-cite, do not re-derive: the selector principle, the epistemic kernel
% and the observation structure are defined in Paper 1. Papers 2-5 cite it.

% SIBLING PAPERS: Paper 1 (Obstruction Certificates under Incomplete Observation), Paper 2 (Exact Probabilistic Sufficiency), Paper 4 (Minimax Dual Certificates), Paper 5 (Exact Belief Computation).
% Cross-cite, do not re-derive: the selector principle, the epistemic kernel
% and the observation structure are defined in Paper 1. Papers 2-5 cite it.

% RESOURCED 2026-09-29: previous assembly (v10) was built from _v9 of this source,
%   11 versions behind. Re-assembled against paper2_computational_certification_v20.tex. Content is a superset:
%   no labels dropped. See VERSION_CURRENCY.md.

\documentclass[10pt,twocolumn]{article}
\usepackage[a4paper,margin=15mm]{geometry}
\usepackage{amsmath,amssymb}
\usepackage{mathrsfs}
\usepackage{amsthm}
\theoremstyle{plain}
\newtheorem{theorem}{Theorem}
\newtheorem{proposition}[theorem]{Proposition}
\newtheorem{corollary}[theorem]{Corollary}
\theoremstyle{definition}
\newtheorem{definition}[theorem]{Definition}
\theoremstyle{remark}
\newtheorem{remark}[theorem]{Remark}
\usepackage{graphicx}
\usepackage{booktabs,array,calc}
\usepackage[font=small,labelfont=bf]{caption}
\usepackage[colorlinks=true]{hyperref}
\providecommand{\tightlist}{\setlength{\itemsep}{0pt}\setlength{\parskip}{0pt}}
\renewcommand{\thesubsection}{\arabic{subsection}}
\makeatletter
\renewcommand{\@seccntformat}[1]{\csname the#1\endcsname.\quad}
\makeatother
\emergencystretch=3em
\allowdisplaybreaks
\newtheorem{conjecture}{Conjecture}
\newtheorem{lemma}{Lemma}
\usepackage{etoolbox}


\begin{document}
\title{Obstruction certificates are small: computable bounds for nonviability and the measure dual that explains why}
\begin{abstract}
\noindent\textbf{The general claim.} A certificate that \emph{no} policy can keep a system within
its constraints is only useful if it is a small object --- finite, attainable by standard solvers,
and independently checkable. This paper establishes that it is, in two independent pipelines, and
identifies what sets its size: \emph{obstruction certificates are small objects whose size is set
by the decision variables, not by the uncertainty.} In the linear-programming pipeline a strictly
positive certified margin is a finite obstruction certificate and, whenever a positive certificate
exists, one can be chosen with at most $r+1$ safety rows --- $r$ an information--time rank, not an
input dimension. In the measure-dual pipeline the certificate is an adversarial probability measure
supported on at most $k+1$ points, $k$ the \emph{control} dimension, and that bound is tight.
Neither the state dimension, nor the cardinality of the disturbance set, nor the resolution of any
discretisation appears in either bound. \emph{The complexity of certifying nonviability is governed
by how many independent decisions the agent has, not by how large the world it operates in is.}
That is what makes these certificates usable rather than merely correct.

\noindent\textbf{Why two parts.} The two pipelines establish the same finding for different objects
and neither contains the other. Part I gives the \emph{computational} pipeline: an outer moment
approximation of all measurable information-adapted controls combined with an inner adversarial
approximation yields a finite linear program whose optimum is a certified lower bound on the
continuous-time safety value and whose error inflation is a certified upper bound; every
dual-feasible solution yields a finite lower-bound certificate covering every measurable policy.
Part II gives the \emph{duality} pipeline: it asks whether the Farkas multiplier that the
polyhedral case happens to carry is an accident of that case or the shadow of something general,
and shows it is the latter --- the obstruction is not a failed search but the existence of a
measure witness. Part I says the certificate can be computed; Part II says what it \emph{is}. The
recurrence of the size bound across them is a result about obstruction certificates, not a
coincidence between two papers, and merging is what makes it visible.

\noindent\textbf{What is found.} Part I: along every admissible refinement sequence whose mesh and
stored-scenario widths vanish, nonviability is eventually detected; and the rank result shows the
bound is not an artefact of the construction --- for every $q$ there is a continuous-time system
with \emph{scalar} input and $q+1$ indistinguishable modes in which every proper subbelief is
viable, the full belief is not, and every certificate must involve all $q+1$ modes. So the
dimension that counts is the information--time product, and the Helly-type bound that would follow
from input dimension is \emph{false} for a common blind control function. Part II: for
control-affine dynamics with compact convex control set, the common safe-action set is empty
precisely when an adversarial measure drives expected inward drift strictly negative against every
control, with the $k+1$ bound tight; the polyhedral case recovers the $\ell_1$-normalised Farkas
multiplier exactly, and the singleton case recovers the Isaacs drift with the adversarial
\emph{distribution} load-bearing. A two-action instance exhibits a strict minimax gap --- without
convexity the equivalence fails and no measure certifies.

\noindent\textbf{Structure.} Part I is the computational certification pipeline and Part II the
measure dual. Each is self-contained and reproduced in full from its companion paper without
condensation, opening with that paper's own abstract. A cross-part conclusion follows, stating the
shared size principle, the two recoveries, and where the equivalence fails.
\end{abstract}
\maketitle

\tableofcontents

\newpage

\part*{Part I --- Computational certification: a finite linear program covering every measurable policy}
\label{part:computational}
\addcontentsline{toc}{part}{Part I --- Computational certification}

\maketitle

\begin{abstract}
Nonviability certificates are useful computationally only if they are finite,
attainable by standard solvers, and independently checkable. This paper
develops such a pipeline for linear systems under exogenous partial
observation of a hidden mode. An outer moment approximation of all measurable
information-adapted controls, combined with an inner adversarial approximation (finitely stored
admissible labels and scenarios), yields a finite linear program whose
optimum is a certified lower bound on the continuous-time safety value and
whose error inflation is a certified upper bound
(\(\rho_{\mathcal{H},G} \le J_{\mathcal{I}}(T) \le \rho_{\mathcal{H},G} +
\bar{e} + \delta_\beta + L h_t\), \(\bar{e}\) the program's own verified
moment-error bound, \(R_U V_U h_u\) in the uniform choice); every dual-feasible solution yields a finite lower-bound certificate
covering every measurable policy, and a strictly positive certified margin
is a finite obstruction certificate; along every admissible refinement
sequence whose mesh and stored-scenario/enclosure widths vanish,
nonviability is eventually detected, and whenever a positive certificate exists one can be
chosen with at most \(r + 1\) safety labels at rank \(r\) --- the relevant
dimension being information--time rank, not instantaneous input dimension,
and a matching construction with scalar input produces obstructions whose
minimal certificates require arbitrarily many witnesses, saturating the
\(r + 1\) bound. The certificates account for post-observation
recourse: on the three-branch delayed-observation instance every pair of compatible
states is viable (explicit pairwise policies are carried exactly) and a
common control keeps every branch safe throughout the blind window, yet nonviability is certified, with the exact
delay threshold \(\tau^{*} = 7/50\); the certificate decodes into an
exact shared-authority delay hierarchy --- the triple is obstructed past
\(\tau = 7/50\), the pairs past \(\approx 0.245\) and \(\approx
0.397\), singletons never --- and the printed blind-window controls are
exactly the optimal shared-braking programs. The certificate is
prescriptive beyond the verdict: advancing observation buys margin at rate
one, dilating post-revelation braking authority buys half a unit of delay
tolerance per unit, and shared-window authority and the velocity budget buy
exactly nothing; belief-cell rows extend the sandwich to partial and noisy
observation with margins radius-bounded on both the witness and the
certificate side. The finite program is solved by a
standard solver (HiGHS) across a five-point mesh refinement; optimality is
certified in exact rational arithmetic by independent primal and dual
witnesses, the solver agreeing with the exact values to at most
\(4.2 \times 10^{-17}\), and the refinement obeys the closed law
\(\rho(h) = 3/50 - Th/4\) (an unrestricted level, nonnegative on this
family for \(h \le \tfrac15\)). A finite-side reference library,
\texttt{viacert}, packages the exact finite-horizon recursion and the
fibre-partition criteria behind a pluggable finite-system interface, with
twelve audited identities in its self-test; the bridge's continuous
programs live in the campaign scripts, not in the library. All exact verification paths for the reported rational benchmark run in
integer and rational arithmetic; solver outputs are numerical
confirmations, never proof objects.
\end{abstract}

\textbf{Keywords:} viability; incomplete observation; obstruction
certificates; linear programming; delayed observation; exact
arithmetic; certified computation

\subsection{Introduction}\label{lp-introduction}

Viability theory under full observation has a mature numerical apparatus:
kernel approximations, set-valued meshes, and convergence results
(Aubin, 1991; Cardaliaguet, Quincampoix, and Saint-Pierre, 1999). Under
incomplete observation the sufficiency side is carried by output-feedback
conditions and the estimation-tube program (Veliov, 1993), while the
necessity side --- certificates whose firing \emph{proves} nonviability of an
information state (Abaee, 2026, An obstruction calculus) --- has lacked a computational bridge. The
difficulty is the quantifier order: the adversary chooses the mode, the
initial state inside its cell, and a measurable disturbance with knowledge of
the announced policy, so naive discretization of controls proves at most the
infeasibility of a held-control subclass, not of the policy class itself.

This paper builds the bridge and the tooling around it. The main result
(Theorem~\ref{lp-thm:bridge}) approximates the admissible controls from outside
by their block moments --- a necessary relaxation of continuous-time safety
that restricts nothing --- and the adversary from inside by explicitly stored
scenarios, yielding a finite linear program whose optimal value sandwiches
the exact minimax safety violation. Dual solutions of the program are finite,
policy-independent obstruction certificates; mesh refinement eventually
certifies --- on a nonviable system sufficiently fine meshes produce a
positive certificate, though positivity is not itself monotone under
refinement (Theorem~\ref{lp-thm:bridge}(iii)); and their size is governed by
the rank of the
safety-row matrix in the information--time product, which no bound on the
instantaneous input dimension controls (Proposition~\ref{lp-prop:rank}).

The pipeline is demonstrated end to end on a three-branch
delayed-observation instance (the running instance of this paper;
Section~\ref{lp-instance}): delay threshold, Farkas certificate, and a
five-point mesh campaign in which the HiGHS solver (Huangfu and Hall, 2018)
agrees with the exact rational optimum at every mesh to floating-point
tolerance, with independent
primal--dual witnesses closing the duality gap exactly. The finite-side
machinery --- the exact backward recursion and the fibre-partition criteria
--- is packaged as the reference library \texttt{viacert} with a pluggable
finite-system interface and a self-test of twelve audited identities; the
library supplies the enumerate-and-check layer of the certification
workflow, and it is deliberately not a reimplementation of the bridge's
continuous programs (Section~\ref{lp-library}).

The contributions are: the continuous-to-finite sandwich theorem with sound,
refinement-complete, sparsity-bounded certificates (Section~\ref{lp-bridge});
the recourse-aware obstruction reading and the exact three-branch instance
with its certified campaign (Section~\ref{lp-instance}); the
information--time rank separation (Section~\ref{lp-rank}); the refinement-erosion
identity, the oracle-recourse extension, and the belief-cell certificates
for partial and noisy observation (Section~\ref{lp-erosion}); the
certificate's prescriptive decode into optimal shared braking authorities,
the induced pair-delay hierarchy, and the exact redesign sensitivities
(Section~\ref{lp-ladder}); the
complexity and residual-bounded verification framework
(Section~\ref{lp-computation}); and the reference library
(Section~\ref{lp-library}).

\subsection{Model, information, and quantifiers}\label{lp-model}

\emph{Dynamics.} Let the hidden mode be
\(\theta \in \Theta := \{1,\dots,M\}\). Conditional on \(\theta\), the
continuous state satisfies
\[
\dot x(t) = A_\theta(t)x(t) + g_\theta(t) + B_\theta(t)u(t) + E_\theta(t)d(t),
\;\; t \in [0,T],
\]
under the standing assumptions: \(A_\theta, g_\theta, B_\theta, E_\theta\)
bounded and piecewise continuous, with bounded-variation bounds on
\(B_\theta, E_\theta\) --- the operative assumption below is the induced
bound \(\mathrm{TV}(k_a) \le V_U\) on the adjoint kernels, assumed
directly in Theorem~\ref{lp-thm:bridge} rather than derived from the
primitive coefficients here; the initial set in mode \(\theta\) a nonempty compact
polytope \(P_\theta \subset \mathbb{R}^n\); the input and disturbance sets
nonempty compact polytopes \(U = \{v : Fv \le f\} \subset \mathbb{R}^m\) and
\(D_\theta \subset \mathbb{R}^p\); safety the closed polyhedron
\(\mathcal{V} = \{x : c_i^\top x \le b_i,\ i = 1,\dots,s\}\); and the initial
information state
\(B_0 = \bigcup_{\theta} \{\theta\} \times P_\theta\) initially safe,
meaning \(P_\theta \subseteq \mathcal{V}\) for every \(\theta\) (if some
\(P_\theta\) leaves \(\mathcal{V}\) at time zero, a time-zero certificate
suffices and the analysis is vacuous). The dynamics admit unique
Carath\'eodory solutions for every admissible measurable signal pair.

\emph{Information.} At known times
\(0 = \tau_0 < \tau_1 < \dots < \tau_R = T\) the controller receives
exogenous information about \(\theta\): refining history partitions
\(\mathcal{P}_0 \preceq \mathcal{P}_1 \preceq \dots \preceq \mathcal{P}_{R-1}\)
of \(\Theta\) --- the order is by refinement: for \(r' > r\) every cell of
\(\mathcal{P}_{r'}\) lies inside a cell of \(\mathcal{P}_r\) (the last
partition is \(\mathcal{P}_{R-1}\) because \(\tau_R = T\) ends the last
interval) --- the observation history during
\(I_r = [\tau_r, \tau_{r+1})\) identifying the cell \(C_r(\theta) \in
\mathcal{P}_r\). These are history partitions, not necessarily partitions
generated by the latest observation alone. Observations do not depend on
\(x, u, d\); active sensing and endogenous observation branching are excluded
from the exact theorem (Section~\ref{lp-erosion} gives a sound extension).

\emph{Policies and quantifiers.} Policies are deterministic, measurable,
nonanticipative, and adapted to the observation history; every such policy is
represented exactly by signals \(u_{r,C} \in L^\infty(I_r; U)\),
\(C \in \mathcal{P}_r\) --- the realized control in mode \(\theta\) being
\(u^{\pi}_\theta(t) = u_{r, C_r(\theta)}(t)\) on \(I_r\). The robust quantifier order is
\[
\exists \pi \in \Pi_{\mathcal{I}}\ \forall \theta\ \forall x_0 \in P_\theta\
\forall d \in L^\infty([0,T]; D_\theta)\ \forall t \in [0,T],
\]
the adversary selecting mode, initial state, and measurable disturbance with
knowledge of the announced policy; randomized policies are not included. Throughout, a belief is a nonempty set \(B \subseteq \Theta \times
\mathbb{R}^n\) of mode-state possibilities, not a probability
distribution. The information-constrained viability kernel is
\(\mathcal{K}_{\mathcal{I}}^T = \{ B : \exists \pi\ \forall \theta, x_0,
d, t\ \ x^{\pi}_{\theta, x_0, d}(t) \in \mathcal{V} \}\), a set of
beliefs; the audited initial state \(B_0\) is viable exactly when
\(B_0 \in \mathcal{K}_{\mathcal{I}}^T\). Writing
\(\mathrm{Lift}(K^T_{\mathrm{full}}) = \{ B \subseteq \Theta \times
\mathbb{R}^n : \varnothing \ne B \subseteq K^T_{\mathrm{full}} \}\) for
the belief lifting of the full-information kernel (an ambient
\(\Theta \times \mathbb{R}^n\)-valued family), \(\mathcal{K}_{\mathcal{I}}^T \subseteq
\mathrm{Lift}(K^T_{\mathrm{full}})\), and the inclusion is strict in general:
the instance of Section~\ref{lp-instance} has every two-state subbelief
viable --- with the pairwise policies computed exactly there --- while the
three-state belief is not.

\medskip
\noindent\textbf{Shared symbols across the companion papers.} The
companion papers share this vocabulary with paper-local meanings; the
letters below carry meanings here that differ there:
\begin{center}\small
\begin{tabular}{@{}p{0.42\linewidth}p{0.52\linewidth}@{}}
\toprule
\textbf{Here} & \textbf{In the companions} \\
\midrule
\(\Gamma_{\mathcal{I}}(\lambda, a)\): the finite obstruction certificate margin (1) & the recourse certificate \(\Gamma_{h}\) (obstruction calculus; minimax duals); the alpha-vector witness sets \(\Gamma^{\Pi}_{k}\) (sufficiency) \\
\(\lambda_{\ell}\): the certificate weights & the Farkas multiplier and observer decay rate (obstruction calculus); the Farkas and dual weights (worked systems) \\
\(\mu\): the adversarial probability measure & Farkas multipliers (obstruction calculus); the adversarial measures of the minimax companion \\
\(K^{T}_{\mathrm{full}}\): the full-information viability kernel & the safe set and kernel \(K\) (obstruction calculus); the window depth and contact set \(\mathcal{K}(B)\) (minimax duals); stock carrying capacities (forecast ladder) \\
\(\rho\): the mesh relaxation parameter & the distributionally robust mixture parameter \(\rho\) (sufficiency; minimax duals); the net biomass multiplier (applied regime viability) \\
\(\beta_{a}\): the stored per-label bound & the margin \(\beta\) (minimax duals) \\
\(\tau\): the observation time & the worst-case exit time \(\tau^{*}\) (worked systems); the delay times \(\tau\) (obstruction calculus) \\
\(\sigma\): the kernel time variable & the floored exit count \(\sigma^{*}\) (worked systems); a noise scale (forecast ladder) \\
\bottomrule
\end{tabular}
\end{center}


\subsection{Adjoint safety rows and the safety value}\label{lp-rows}

Let \(X_\theta(t,s)\) denote the state-transition matrix. For a label
\(a = (\theta, i, t) \in \mathscr{S} := \Theta \times \{1,\dots,s\} \times
[0,T]\) define the input kernel and margin
\[
k_a(\sigma) = \mathbf{1}_{\{\sigma \le t\}} B_\theta(\sigma)^\top
X_\theta(t,\sigma)^\top c_i,
\]
\begin{align*}
\beta_a = \;& b_i - h_{P_\theta}\!\bigl(X_\theta(t,0)^\top c_i\bigr)
- \int_0^t c_i^\top X_\theta(t,\sigma) g_\theta(\sigma)\, d\sigma\\
& - \int_0^t h_{D_\theta}\!\bigl(E_\theta(\sigma)^\top
X_\theta(t,\sigma)^\top c_i\bigr)\, d\sigma,
\end{align*}
where \(h_Q(v) = \sup_{q \in Q} v^\top q\). For a policy \(\pi\) set
\(\mathcal{F}_a(\pi) = \int_0^T k_a(\sigma)^\top u^\pi_\theta(\sigma)\,
d\sigma - \beta_a\).

\begin{proposition}[robust rows]\label{lp-prop:rows}
\(\pi\) is robustly safe on \([0,T]\) if and only if
\(\mathcal{F}_a(\pi) \le 0\) for every \(a \in \mathscr{S}\).
\end{proposition}

\begin{proof}
The identity is the linear solution formula: the initial-state and
disturbance maximizations separate from the control term because the
control signal in a mode does not depend on \(x_0\) or \(d\), so the
rows are exact. For polytopic \(D_\theta\), a maximizing disturbance
can be chosen measurably by selecting a maximizing vertex at each time,
ties resolved by a fixed ordering.
\end{proof}

Define the minimax safety violation
\(J_{\mathcal{I}}(T) = \min_{\pi \in \Pi_{\mathcal{I}}} \max_{a \in
\mathscr{S}} \mathcal{F}_a(\pi)\).

\begin{proposition}[value and nonviability]\label{lp-prop:value}
The minimum exists: the policy-signal space is a finite product of
weak-* compact sets \(L^\infty(I_r; U)\) (each \(U\) a compact convex
polytope, so the admissible signal sets are weak-* closed and bounded) and
each \(\mathcal{F}_a\) is weak-* continuous; the maximum over labels is
attained on the compact label set \(\mathscr{S}\) because
\(a \mapsto \mathcal{F}_a(\pi)\) is continuous in \(t\) under
assumption (6) below. Consequently
\(B_0 \notin \mathcal{K}_{\mathcal{I}}^T\) if and only if
\(J_{\mathcal{I}}(T) > 0\).
\end{proposition}

\begin{proof}
Existence: the policy-signal space is a finite product of weak-*
compact sets \(L^\infty(I_r; U)\) (each \(U\) a compact convex
polytope, so the admissible signal sets are weak-* closed and bounded)
and each \(\mathcal{F}_a\) is weak-* continuous; the maximum over
labels is attained on the compact label set \(\mathscr{S}\) because
\(a \mapsto \mathcal{F}_a(\pi)\) is continuous in \(t\) under
assumption (6) below. The equivalence is then the definition of the
kernel: a violation-free policy exists exactly when the minimax value
vanishes. The strict positive margin is a consequence of compactness
and closed safety constraints, not an extra nonviability assumption.
\end{proof}

\subsection{The finite obstruction certificate}\label{lp-certificate}

Define \(\phi_U(g) := \min_{v \in U} g^\top v = -h_U(-g)\). Choose finitely
many labels \(a_\ell = (\theta_\ell, i_\ell, t_\ell)\), \(\ell = 1,\dots,q\),
and weights \(\lambda_\ell \ge 0\), \(\sum_\ell \lambda_\ell = 1\). For each
active information cell \(C \in \mathcal{P}_r\), pool the adjoint kernels
\(g_{r,C}(\sigma) = \sum_{\ell : \theta_\ell \in C} \lambda_\ell
k_{a_\ell}(\sigma)\). The certificate margin is
\[
\Gamma_{\mathcal{I}}(\lambda, a)
= \sum_{r=0}^{R-1} \int_{I_r} \sum_{C \in \mathcal{P}_r}
\phi_U\!\bigl(g_{r,C}(\sigma)\bigr)\, d\sigma
- \sum_{\ell=1}^q \lambda_\ell \beta_{a_\ell}.
\tag{1}
\]

\begin{remark}[semantics]\label{lp-rem:semantics}
If \(\Gamma_{\mathcal{I}}(\lambda, a) > 0\), then no admissible
information-adapted policy is safe: for every policy,
\(\sum_\ell \lambda_\ell \mathcal{F}_{a_\ell}(\pi) \ge
\Gamma_{\mathcal{I}}(\lambda, a)\), so at least one stored label satisfies
\(\mathcal{F}_{a_\ell}(\pi) \ge \Gamma_{\mathcal{I}}(\lambda, a) > 0\) ---
an admissible realization violates facet \(i_\ell\) at time \(t_\ell\). The
violating label may depend on the policy; the finite list and its weights do
not. The certificate is not an adversarial tree: it is a finite,
policy-independent set of adversarial tests combined by an adjoint identity.
The pooled bound is tight as a minimum over policies:
\(\min_{\pi} \sum_\ell \lambda_\ell \mathcal{F}_{a_\ell}(\pi) =
\Gamma_{\mathcal{I}}(\lambda, a)\), the pointwise minimizer being
measurable by fixed-order vertex selection in the polytope \(U\).
\end{remark}

\subsection{The continuous-to-finite theorem}\label{lp-bridge}

The exact dual expression (1) is a genuine infinite-dimensional minimax
statement, and we prove it in full rather than invoke it: the proof of
(iv) below (complete version in the Supplementary Material) runs Sion's
minimax theorem over the mixed extension to probability measures on the
label set and transfers the result to finitely supported measures by
weak-* density. The consequential result is a sound finite construction
covering unrestricted measurable controls; classical infinite-dimensional
duality (Anderson and Nash, 1987; Hettich and Kortanek, 1993) is the
surrounding theory, not a load-bearing citation.

\emph{Outer control moments.} Choose a control-time mesh
\(\mathcal{H} = \{J_1,\dots,J_H\}\) containing every sensing time. Introduce
one vector \(v_{j,C} \in U\) for every cell \(C \in \mathcal{P}_{r(j)}\)
active during \(J_j\), where \(r(j)\) is the unique observation interval
containing \(J_j\) (the mesh contains every sensing time, so each block
lies in one \(I_r\)), and let \(Q = \sum_j |\mathcal{P}_{r(j)}|\) be the
number of information--time blocks. For a label \(a\) define the finite row
\(A_{a,j,C} = \mathbf{1}_{\{\theta \in C\}} \int_{J_j} k_a(\sigma)\, d\sigma\),
the block average
\(\bar{k}_{a,\mathcal{H}}(\sigma) = \frac{1}{|J_j|} \int_{J_j} k_a(s)\, ds\)
on \(J_j\), and the moment error
\[
e_{a,\mathcal{H}} = \int_0^T h_U\!\bigl(\bar{k}_{a,\mathcal{H}}(\sigma)
- k_a(\sigma)\bigr)\, d\sigma,
\tag{2}
\]
for which any verified upper bound \(\bar{e}_a \ge e_{a,\mathcal{H}}\) may be
used. For every measurable policy, its block averages belong to \(U\), and
\[
\sum_{j,C} A_{a,j,C}^\top v_{j,C} - \int_0^T k_a(\sigma)^\top
u^\pi_{\theta(a)}(\sigma)\, d\sigma \le \bar{e}_a.
\tag{3}
\]
Inequality (3) is the crucial orientation: finite feasibility is a
necessary condition for continuous-time safety. The moments do not restrict
the real controller to held inputs. For the sampled-data subclass whose policies are
held constant on the blocks, the averaging term vanishes identically
(\(\int_{J_j}(k_a - \bar{k}_{a,\mathcal{H}})\, d\sigma = 0\) on every
block), so such a policy is continuously safe exactly when its block
moments satisfy the unrelaxed rows: an exact reduction precisely when the
row bounds are exact, \(\bar{e}_a = 0\); with the error bounds (2) the
program remains a relaxation even for held policies (in the instance the
held witness policy attains the continuous value \(\tfrac3{50}\) on every
aligned mesh, while (5) returns \(\tfrac3{50} - Th/4\)).

\emph{Inner adversarial approximation.} One may use a right-hand side
\(\beta_a^{+} \ge \beta_a\) obtained from an explicit admissible initial
state and disturbance signal rather than exact worst-case optimization. If
\(0 \le \beta_a^{+} - \beta_a \le \delta_\beta\) uniformly over the labels
used, the approximation is safe for nonviability: it weakens the adversary
and enlarges the feasible controller set. For example, with
\(w_a(\sigma) = \mathbf{1}_{\{\sigma \le t\}}
E_\theta(\sigma)^\top X_\theta(t,\sigma)^\top c_i\), on a disturbance mesh
choose \(d_{a,J} \in \arg\max_{d \in D_\theta} (\int_J w_a)^\top d\),
yielding an explicit piecewise-constant admissible disturbance; if
\(D_\theta\) has radius \(R_D\) about a fixed center and
\(\mathrm{TV}(w_a) \le V_D\), then
\[
0 \le \beta_a^{+} - \beta_a \le R_D V_D h_d,
\tag{4}
\]
the general budget being \(\delta_\beta = \delta_{\mathrm{init}} +
\delta_{\mathrm{dist}} + \delta_{\mathrm{num}}\) for separately accounted
initial-state and numerical error.

\emph{The finite program.} Let \(G \subset [0,T]\) be a finite time grid with \(0, T \in G\) and
every observation time in \(G\); the label set is the product
\(\mathscr{S}_G = \Theta \times \{1,\dots,s\} \times G\), with covering
radius \(h_t = \sup_{t \in [0,T]} \min_{\hat{t} \in G} |t -
\hat{t}|\). Solve
\begin{align*}
\rho_{\mathcal{H},G} = \min_{v, \rho}\ \rho
\quad \text{s.t.} \quad
& A_a v - \beta_a^{+} - \bar{e}_a \le \rho \\
& (a \in \mathscr{S}_G),
\ v_{j,C} \in U.
\tag{5}
\end{align*}
The level \(\rho\) is an unrestricted real variable; negative values record
strict safety slack.

\begin{theorem}[continuous-to-finite]\label{lp-thm:bridge}
Assume:
\begin{itemize}\setlength\itemsep{0em}
\item[(H1)] uniformly over admissible policies, modes, and facets,
\(|\mathcal{F}_{\theta,i,t}(\pi) - \mathcal{F}_{\theta,i,t'}(\pi)| \le
L|t - t'|\) (6);
\item[(H2)] the input kernels satisfy \(\mathrm{TV}(k_a) \le V_U\) (7);
\item[(H3)] the label maps \(a \mapsto k_a\) and \(a \mapsto \beta_a\)
are continuous on the compact label set (\(k_a\) in \(L^1\) norm,
\(\beta_a\) absolutely; the \(L^1\) form is the operative one ---
sup-norm continuity is neither needed nor true here, the step kernels
below jumping at \(\sigma = t\) for every \(t\) --- while \(L^1\)
continuity is implied by bounded coefficients and enjoyed by every
rational instance below);
\item[(H4)] \(U\) has radius \(R_U\) about some fixed center;
\item[(H5)] the time grid has covering radius at most \(h_t\) and the
control mesh maximum length is \(h_u\);
\item[(H6)] the adversarial approximation satisfies the bound above.
\end{itemize}
Then:
\end{theorem}

\noindent\emph{(i) Certified value sandwich.}
\[
\rho_{\mathcal{H},G} \le J_{\mathcal{I}}(T) \le \rho_{\mathcal{H},G}
+ \max_{a \in \mathscr{S}_G} \bar{e}_a + \delta_\beta + L h_t.
\tag{8}
\]
With (4) --- the budget \(\delta_\beta = \delta_{\mathrm{init}} +
\delta_{\mathrm{dist}} + \delta_{\mathrm{num}}\) --- \(J_{\mathcal{I}}(T) -
\rho_{\mathcal{H},G} \le \max_{a} \bar{e}_a + R_D
V_D h_d + \delta_{\mathrm{init}} + \delta_{\mathrm{num}} + L h_t\).
Verified numerical row widths are added to the right side. The upper bound
is realized by held controls, for which the averaging error vanishes
identically, so the bound is charged by the error bounds the program
actually used; the uniform choice \(\bar{e}_a \equiv R_U V_U h_u\) ---
valid by (2) whenever \(\mathrm{TV}(k_a) \le V_U\), since
\(\int_J |k_a - \bar{k}_a| \le h_u\, \mathrm{TV}(k_a|_J)\) and
\(h_U(g) \le R_U \|g\|\) --- recovers the announced
\(R_U V_U h_u\) form as a special case. The \(\max \bar{e}_a\) form is the conservative one: a loose but verified choice inflates \(\rho\)'s relaxation
and must inflate the charged bound with it (a counterexample is carried in
Section~\ref{lp-instance}).

\noindent\emph{(ii) Sound finite certificates.} If \(\rho_{\mathcal{H},G} > 0\), the
program's dual supplies a finite certificate of nonviability of the original
continuous-time system, covering every measurable information-adapted policy.

\noindent\emph{(iii) Completeness under refinement.} If
\(B_0 \notin \mathcal{K}_{\mathcal{I}}^T\), then for every admissible
refinement sequence along which \(h_u, h_d, h_t\) converge to zero and the
stored-scenario widths \(\delta_{\mathrm{init}}, \delta_{\mathrm{num}}\)
do as well --- near-maximizing stored initial states
(Section~\ref{lp-computation}) must be refined with the mesh for completeness
to hold --- the meshes eventually produce a positive
finite certificate. Nonviability is thus semidecidable in this class for
the input model in which the coefficients are exact rational or
interval-enclosed piecewise-polynomial data with certified integration,
total-variation, and uniform time-Lipschitz bounds: the semidecision
procedure enumerates admissible meshes, solves each enclosed rational
program, and halts on the first positive certificate --- for instance, the
dyadic schedule \(h_u = h_d = h_t = T/2^k\) halts with certificate margin
at least \(\gamma/2\) once \(2^{-k}(R_U V_U + R_D V_D + L)T < \gamma/2\)
(with exactly stored scenarios, \(\delta_{\mathrm{init}} =
\delta_{\mathrm{num}} = 0\)),
where \(\gamma = J_{\mathcal{I}}(T) > 0\) is the continuous value.

\noindent\emph{(iv) Exact atomic dual representation.}
\[
J_{\mathcal{I}}(T) = \sup_{q < \infty,\ a_1,\dots,a_q,\
\lambda \in \Delta_q} \Gamma_{\mathcal{I}}(\lambda, a).
\tag{9}
\]
In particular, nonviability is equivalent to the existence of a positive
finite atomic certificate.

\noindent\emph{(v) Sparse witnesses.} Let \(r = \mathrm{rank}\, A\) for the
finite safety-row coefficient matrix of (5) --- the \(A_a\) rows excluding
the \(\rho\)-column and the actuator inequalities; the bound is linear-programming (LP)
basic-solution theory (Chv\'{a}tal, 1983; supplementary Step 6). Whenever \(\rho_{\mathcal{H},G} > 0\), a
positive certificate exists using at most \(r + 1 \le mQ + 1\) safety labels,
and it can retain any prescribed margin strictly below
\(\rho_{\mathcal{H},G}\). If the lifted continuous-time kernels span a
finite-dimensional space of dimension \(r_{\mathcal{I}}\), a continuous-time
obstruction admits one with at most \(r_{\mathcal{I}} + 1\) labels. The relevant
dimension is information--time rank, not \(m\).

\emph{Proof (the complete steps are Steps 1--6 of the supplementary
material, verified by the verification script).} The lower bound holds because every admissible
measurable policy's control bounds are outer-approximated by the block
moments, so the moment program relaxes the controller's feasible set and
its value cannot exceed the continuous safety value; the upper bound is
realized by taking, on each block, the piecewise-constant control attained
by the finite optimizer's block variables, whose exit from the continuous
value is charged exactly by the program's own moment bounds \(\bar{e}_a\)
(\(R_U V_U h_u\) under the uniform choice), the stored-scenario error
\(\delta_\beta\), and the time-grid interpolation error \(L h_t\). The asymmetry at its core is
worth stating plainly: averaging arbitrary controls proves the lower bound;
implementing piecewise-constant controls proves the upper bound. Confusing
these two directions is what makes na\"ive infeasibility-by-discretization
arguments unsound.

\emph{Proof of (iv) (exact atomic dual representation).} The policy space
\(\Pi_{\mathcal{I}} = \prod_{r,C}\{u \in L^\infty(I_r; U)\}\) is
weak-* compact and convex (Banach--Alaoglu: each factor is weak-* closed
and bounded, \(U\) being a compact convex polytope), and the label set
\(\mathscr{S}\) is compact. Under the label-continuity hypothesis,
\((a, \pi) \mapsto \mathcal{F}_a(\pi)\) is jointly continuous
(\(\pi\)-continuity is the \(L^1\)-pairing; \(a\)-continuity is
uniform in \(\pi\) by the \(L^1\) label continuity and boundedness of
the \(U\)-values: \(|\mathcal{F}_a(\pi) - \mathcal{F}_{a'}(\pi)| \le
M_U \|k_a - k_{a'}\|_{L^1} + |\beta_a - \beta_{a'}|\) with
\(M_U := \sup_{v \in U} \|v\|\) (H\"older; the controls are
\(M_U\)-bounded), the equicontinuity
Arzel\`a--Ascoli consumes). Let \(\mathcal{P}(\mathscr{S})\) carry the weak-*
topology: by the Riesz representation it is the compact convex metrizable
set of Borel probability measures. Since the integral of a continuous
function over a probability measure lies within its min--max range and
the Dirac measures attain both ends,
\(\max_a \mathcal{F}_a(\pi) = \max_\mu \int
\mathcal{F}_a(\pi)\,d\mu\), and the bilinear
\(\int \mathcal{F}_a(\pi)\,d\mu\) --- affine, hence simultaneously
quasi-convex and quasi-concave, and continuous in each variable --- meets
the hypotheses of Sion's minimax theorem, giving
\[
J_{\mathcal{I}}(T) \;=\; \max_{\mu \in \mathcal{P}(\mathscr{S})}
\; \min_{\pi} \int \mathcal{F}_a(\pi)\,d\mu,
\]
the outer maximum attained. The value function
\(\mu \mapsto \min_\pi \int \mathcal{F}_a(\pi)\,d\mu\) is weak-*
continuous: the family \(\{a \mapsto \mathcal{F}_a(\pi)\}_{\pi}\)
is uniformly bounded and uniformly equicontinuous, so by
Arzel\`a--Ascoli weak-* convergence of measures makes the pairings
converge uniformly over \(\pi\). Finitely supported measures are weak-*
dense in \(\mathcal{P}(\mathscr{S})\), so this continuous value
function has the same supremum over finitely supported measures as over
all of \(\mathcal{P}(\mathscr{S})\) --- namely \(J_{\mathcal{I}}(T)\).
For \(\mu = \sum_\ell \lambda_\ell \delta_{a_\ell}\) the inner
minimum decomposes over information cells, and the cellwise infimum of an
\(L^1\)-pairing over \(U\)-valued signals equals the pointwise
\(\phi_U\) integral by measurable vertex selection (fixed vertex order
in the polytope), which is exactly the representation
\(\Gamma_{\mathcal{I}}(\lambda, a)\) of (1). Nonviability
\(J_{\mathcal{I}}(T) > 0\) therefore yields a finitely supported
\(\mu\) with \(\Gamma > J_{\mathcal{I}}(T)/2 > 0\); conversely
\(\Gamma > 0\) lower-bounds \(J_{\mathcal{I}}(T)\) because
\(\max_a \ge \) the \(\lambda\)-weighted average. The complete
seven-step verification of the remaining hypotheses is Step 5b of the
Supplementary Material; the label-continuity input is the \(L^1\)
estimate above, verified in Section~\ref{lp-instance} for the audited
kernels.

\noindent\emph{Corollary}[complete dual calculus].
\label{lp-cor:dual}
Under the hypotheses of Theorem~\ref{lp-thm:bridge}: (a) every dual-feasible
pair \((\lambda, a)\) satisfies
\(\Gamma_{\mathcal{I}}(\lambda, a) \le J_{\mathcal{I}}(T) \le
\max_a \mathcal{F}_a(\pi)\) for every policy --- the lower bound is
Remark~\ref{lp-rem:semantics}, the upper is the weighted-average bound of
the proof; (b) \(J_{\mathcal{I}}(T) = \sup_{q, a, \lambda \in
\Delta_q} \Gamma_{\mathcal{I}}(\lambda, a)\): the supremum over all
measures is attained (Sion), and finitely supported measures approach its
value arbitrarily finely --- attainment by an atomic maximizer is an
instance-level fact (in the instance of Section~\ref{lp-instance} three
atoms attain it), not a general claim; (c) \(B_0 \notin
\mathcal{K}_{\mathcal{I}}^T\) if and only if some dual-feasible pair has
\(\Gamma_{\mathcal{I}}(\lambda, a) > 0\); every such pair obstructs
every measurable information-adapted policy, and every obstruction arises
this way. The duals therefore form the complete certificate system for
the class, positivity being the only decision.

\subsection{A three-branch delayed-observation instance carried exactly}\label{lp-instance}

\emph{Plant and exact delay threshold.} Consider the planar double integrator
\(\dot p = v\), \(\dot v = u\), with unit vectors \(n_1 = (1,0)\),
\(n_2 = (-\tfrac35, \tfrac45)\), \(n_3 = (-\tfrac35, -\tfrac45)\), the
hexagonal input set
\(U = \mathrm{conv}\{\pm n_1, \pm n_2, \pm n_3\}\), equivalently
\(\{\nu : |u_2| \le \tfrac45,\ |2u_1 + u_2| \le 2,\ |2u_1 - u_2| \le 2\}\),
safety \(n_i^\top p \le 2\) (\(i = 1,2,3\)) and
\(|v_1|, |v_2| \le \tfrac65\) --- seven facets in all, \(s = 7\) (three
position, four velocity) --- three possible initial states
\(p_j(0) = \tfrac{34}{25} n_j\), \(v_j(0) = n_j\), no disturbance, and the
mode unobserved until time \(\tau\), when it is revealed exactly.

If \(j\) is known, \(u = -n_j\) for one second and then \(u = 0\) keeps the
maximum critical position at \(\tfrac{34}{25} + \tfrac12 =
\tfrac{93}{50} = 1.86 < 2\): all three physical states are indefinitely
viable under full information (throughout the braking leg
\(\|v(t)\|_\infty \le 1 \le \tfrac65\); the noncritical facets keep
\(n_i^\top p < 0\) along the \(n_j\)-ray; \(-n_j \in U\)). Take three witness labels, each using its own position facet at
\(t_* = \tau + 1\), with kernels
\(k_j(s) = \mathbf{1}_{\{s \le t_*\}} (t_* - s) n_j\) and
\(\beta_j = \tfrac{16}{25} - t_*\), weighted by \(\lambda_1 = \tfrac38\),
\(\lambda_2 = \lambda_3 = \tfrac{5}{16}\), which satisfy \(\sum_j \lambda_j = 1\) and \(\sum_j \lambda_j n_j =
0\). Before observation the pooled kernel is
identically zero; after observation each mode has its own control and
\(\phi_U(\lambda_j (t_* - s) n_j) = -\lambda_j (t_* - s)\). Hence
\[
\Gamma(\tau) = -\int_\tau^{\tau+1} (\tau + 1 - s)\, ds
- \Bigl(\tfrac{16}{25} - \tau - 1\Bigr) = \tau - \frac{7}{50},
\tag{10}
\]
so \(\tau > \tfrac{7}{50}\) implies
\(B_0 \notin \mathcal{K}_{\mathcal{I}}\). Conversely, \(u = 0\) before
observation and \(u = -n_j\) for one second afterward keeps the maximum
critical position at \(\tfrac{34}{25} + \tau + \tfrac12\), safe exactly when
\(\tau \le 0.14\). Thus \(\tau_{\max} = \tfrac{7}{50} = 0.14\) is exact (equality viable:
the safety set is closed), and for horizons \(T \ge \tau + 1\) the
upper-bound policy --- full braking from \(\tau\) on, stopping at
\(\tau + 1\) from \(p = (\tfrac{34}{25} + \tau)n_j\) --- shows every row
value at most \(\tau - \tfrac{7}{50}\), so
\(J_{\mathcal{I}}(T) = \tau - \tfrac{7}{50}\).

\emph{Pairs are viable; the blind-window exit test misses.} At
\(\tau = \tfrac15\), the common control \(u = 0\) keeps all three branches
safe throughout the blind window:
\(n_j^\top p_j(0.2) = \tfrac{39}{25} = 1.56 < 2\). The pair claims are
carried exactly. For \(\{1,2\}\) (and \(\{1,3\}\) by the reflection
\(v_2 \mapsto -v_2\)): brake \emph{through} the blind window with the
common control \(u = -(n_1 + n_2) = (-\tfrac25, -\tfrac45) \in U\) (both
branches decelerate along their normals at \(\tfrac25\)), then brake each
branch along its own normal, \(u = -n_j\) from \(\tau\) to
\(n_j\)-rest --- the critical position peaks at
\(\tfrac{12345}{6250} = 1.9752 < 2\) on both branches, slack
\(\tfrac{31}{1250}\). For \(\{2,3\}\): hold \(u = n_1\) through the
blind window (both branches decelerate at \(\tfrac35\)), then brake each
branch along its own normal, \(u = -n_j\) from \(\tau\) --- the critical
position peaks at \(\tfrac{2419}{1250} = 1.9352\), slack
\(\tfrac{81}{1250}\). (Two tempting variants fail: holding \(u = 0\)
through the blind window in \(\{1,2\}\) peaks at
\(\tfrac{2719}{1250} = 2.175 > 2\) with the printed mid-brake, and at
\(\tfrac{103}{50}\) even with full braking afterward; inserting a
\(-\tfrac35 n_j\) leg after \(\tau\) in \(\{2,3\}\) peaks at
\(\tfrac{2501}{1250} = 2.0008 > 2\) at \(t = \tfrac{29}{25}\).) All
seven facets hold at every time on every in-belief branch of every pair
policy on \([0, \tfrac65]\): every velocity value along \(\{1,2\}\)
satisfies \(\|v\|_\infty \le 1\) and along \(\{2,3\}\)
\(\|v\|_\infty \le \tfrac45\), both strict against
\(\tfrac65\), and the worst slacks are the peak rows
\(\tfrac{31}{1250}\) and \(\tfrac{81}{1250}\); braking to
\(n_j\)-rest leaves a harmless transverse velocity of \(\tfrac4{25}\)
(its drift only decreases the noncritical coordinates). The blind control
\(u = n_1\) does drive the off-belief branch 1 to \(|v_1| = \tfrac65\)
exactly at \(\tau\) --- irrelevant to the pair claim (branch 1 is not in
the belief), and indeed that control could not keep branch 1 safe: its
peak would be \(\tfrac{23}{10} = 2.3 > 2\). No exit occurs
before information arrives; nevertheless \(\Gamma(0.2) = 0.06 > 0\) --- the best
post-observation recourse cannot recover all branches. The obstruction is
not that information arrives after unavoidable physical exit; it is that
information arrives after some indistinguishable branch has necessarily
accumulated more braking liability than the remaining control authority can
discharge.

\begin{figure}[t]
\centering
\includegraphics[width=0.94\linewidth]{figs_comp2/fig_trajectories.pdf}
\caption{The audited policies on the critical coordinate
\(n_j^{\top} p_j(t)\) at \(\tau = \tfrac15\), every peak the exact
rational of the text: the full-information singleton peaks at
\(\tfrac{93}{50}\); the pair policies peak below the floor
(\(\tfrac{12345}{6250}\), \(\tfrac{2419}{1250}\), slacks
\(\tfrac{31}{1250}\) and \(\tfrac{81}{1250}\)); the two falsified
variants cross it (\(\tfrac{103}{50}\); \(\tfrac{2501}{1250}\) at
\(t = \tfrac{29}{25}\)). Generated by the assertion script
\texttt{paper2\_comp\_certification\_figures\_v1.py}.}
\label{lp-fig:trajectories}
\end{figure}

\emph{The finite program and its exact witnesses.} Set
\(\tau = \tfrac15\), \(T = \tfrac65\), \(h = \tfrac1{10}\): two common
pre-observation control blocks and ten post-observation blocks for each of
three modes, so \(Q = 32\), \(mQ = 64\). For the endpoint position rows the
exact moment error is \(e = 12 \cdot h^2/4 = \tfrac{3}{100}\) (the kernel is
affine on each of the \(T/h = 12\) mesh intervals and
\(h_U(\alpha n_j) = |\alpha|\)); the endpoint right-hand side is
\(\beta = \tfrac{16}{25} - \tfrac65 = -\tfrac{14}{25}\), relaxed to
\(\beta + e = -\tfrac{53}{100}\). The exact continuous certificate has
margin \(\tfrac{3}{50} = 0.06\); the finite moment certificate has margin
\(\tfrac{3}{100} = 0.03\).

The dual witness is exact. Order the input facets by
\(F = [(0,1);(0,-1);(2,1);(2,-1);(-2,1);(-2,-1)]\),
\(f = (\tfrac45, \tfrac45, 2, 2, 2, 2)\), and define
\(\eta_1 = (0,0,0,0,\tfrac14,\tfrac14)\),
\(\eta_2 = (0,\tfrac12,0,\tfrac3{10},0,0)\),
\(\eta_3 = (\tfrac12,0,\tfrac3{10},0,0,0)\), satisfying
\(F^\top \eta_j = -n_j\) and \(f^\top \eta_j = 1\). For a post-observation
interval \(J_k\) put \(a_k = \int_{J_k} (T - s)\, ds\) and
\(\mu_{j,k} = \lambda_j a_k \eta_j\). Before observation, stationarity holds
because \(\sum_j \lambda_j n_j = 0\); after observation, because
\(\lambda_j a_k n_j + F^\top \mu_{j,k} = 0\). Since
\(\sum_k a_k = \tfrac12\), the Farkas contradiction is
\(\lambda^\top(\beta + e) + \mu^\top f_{\mathrm{blk}} = -\tfrac{53}{100} +
\tfrac12 = -\tfrac{3}{100} < 0\) --- an exact rational certificate. The
primal witness is equally exact: hold the pre-observation blocks at \(0\) and
the post-observation blocks of mode \(j\) at \(-n_j\); the endpoint row then
reads \(-\tfrac12 + \tfrac{14}{25} - \tfrac{3}{100} = \tfrac{3}{100}\), and
every other row lies at or below the same value. Three structural facts
make this transparent: the velocity and non-critical position rows are
continuously safe at the witness (\(\|v_\theta(t)\|_\infty \le 1\)
against the \(\tfrac65\) boundary; \(n_i^\top n_j < 0\) for
\(i \ne j\) with the witness trajectory on the \(n_j\) ray), so their
row values \(\mathcal F_a(\hat\pi) - \bar e_a\) are nonpositive; the
endpoint critical rows attain \(\rho = \tfrac3{100}\) with equality; and
the remaining critical rows interpolate monotonically between the two
(the braking profile \(\zeta(t) = \tfrac{39}{25} + (t - \tau) -
\tfrac{(t-\tau)^2}{2}\) increases to its endpoint maximum) --- these are
verified row-by-row, with the campaign script checking all \(465\) rows
exactly at this mesh and archiving the slacks
(Section~\ref{lp-methods}).

\emph{Instance arithmetic, collected.} The exact rational chain of this
subsection is: the pooling weights
\(\lambda = (\tfrac38, \tfrac5{16}, \tfrac5{16})\) with
\(\sum_j \lambda_j = 1\) and \(\sum_j \lambda_j n_j = 0\); the
witness margins \(\beta_j = \tfrac{16}{25} - t_*\) and the value
\(\Gamma(\tau) = \tau - \tfrac7{50}\) in (10), so
\(\tau_{\max} = \tfrac7{50}\) with equality viable; the moment error
\(e = Th/4\), giving at \(h = \tfrac1{10}\) the relaxed right-hand
side \(\beta + e = -\tfrac{53}{100}\), the dual Farkas value
\(-\tfrac{53}{100} + \tfrac12 = -\tfrac3{100} < 0\), and the primal
endpoint row \( -\tfrac12 + \tfrac{14}{25} - \tfrac3{100} =
+\tfrac3{100}\); and the law \(\rho(h) = \tfrac3{50} - Th/4\) on every uniform aligned
mesh (the level unrestricted; negative on coarser non-uniform aligned
meshes). Every identity is an
exact rational statement, machine-verified in
Section~\ref{lp-methods}.

\emph{Solver-assisted, exactly certified mesh study.} For each mesh the
program (5) is built
at its claimed dimensions and solved by HiGHS; optimality is then certified
in exact rational arithmetic by the witness pair above (primal feasibility of
every row; dual stationarity and dual objective \(-\rho(h)\)), so the duality
gap closes exactly and the solver's returned value is a computational
confirmation of a proved optimum, not an observation.
Table~\ref{lp-tab:mesh} reports the campaign: exact values
\(\rho(h) = \tfrac{3}{50} - Th/4\) (nonnegative on this family), solver values, iteration
counts, and wall-clock times as measured on the verification machine.

\begin{table*}[t]
\centering
\footnotesize
\begin{tabular}{@{}r|r|r|r|r|r|r|r@{}}
\toprule
\(h\) & blocks \(Q\) & variables \(mQ{+}1\) & rows \(Ms(N_t{+}1)+p_UQ\) & exact \(\rho(h)\) & solver \(\rho\) & iterations & time (s) \\
\midrule
0.20 & 16 & 33  & 243  & \(0\) & 0.000000000 & 36   & 0.024 \\
0.10 & 32 & 65  & 465  & \(\tfrac{3}{100} = 0.030\) & 0.030000000 & 78   & 0.014 \\
0.05 & 64 & 129 & 909  & \(\tfrac{9}{200} = 0.045\) & 0.045000000 & 141  & 0.043 \\
0.02 & 160 & 321 & 2241 & \(\tfrac{27}{500} = 0.054\) & 0.054000000 & 501  & 0.252 \\
0.01 & 320 & 641 & 4461 & \(\tfrac{57}{1000} = 0.057\) & 0.057000000 & 1062 & 1.015 \\
\bottomrule
\end{tabular}
\caption{Solver-assisted, exactly certified mesh study on the three-branch
instance (\(\tau = \tfrac15\), \(T = \tfrac65\)): program dimensions,
exact optimal values, returned values (agreement to at most
\(4.2 \times 10^{-17}\)), simplex iterations, and wall-clock time (single
runs, scipy HiGHS dual simplex, on the verification machine; iteration
counts and times may vary with solver version and settings; ``rows'' counts
LP inequalities --- safety plus actuator, with \(s = 7\), \(p_U = 6\),
\(N_t = T/h\)). Optimality of each exact value is certified by the exact
primal--dual witness pair, independent of the solver. At \(h = 0.2\) the
witness is an exact optimality witness of value zero; at the four finer
meshes it is a positive obstruction certificate. The exact continuous value
is \(0.06\); the three-label support of each positive certificate is an
instance-specific fact of this symmetric campaign, distinct from the
general \(r + 1\) bound.}
\label{lp-tab:mesh}
\end{table*}

On the audited family of uniform meshes \(h \in \{\tfrac15,
\tfrac1{10}, \tfrac1{20}, \tfrac1{50}, \tfrac1{100}\}\) (each dividing
both \(\tau\) and \(T\)), the values obey the exact law \(\rho(h) =
\tfrac{3}{50} - Th/4\) for \(h < \tfrac15\) and vanish at \(h =
\tfrac15\): the whole refinement gap is the control-moment error \(Th/4\)
of the endpoint rows: the finite value lies below the continuous margin by
exactly that error and converges to it from below as \(h \downarrow 0\)
(\(\rho \le J\) always --- the program is a relaxation; Table~\ref{lp-tab:mesh}
itself shows \(0.03 < 0.06\)). The \(\max\{0,\,\cdot\}\) in the
table's law is inactive on this family: for \(h \le \tfrac15\) the raw
value \(3/50 - Th/4\) is already nonnegative, vanishing at
\(h = \tfrac15\); the level is unrestricted, and coarser non-uniform
meshes record negative values (below).
For the record, the instance's hypothesis constants are
\(R_U = 1\) (the hexagon's circumradius), \(V_U = \tau + 1 = \tfrac65\)
(the endpoint kernels are unit-slope ramps of height \(t^{*}\)),
\(\delta_\beta = 0\) (\(\beta\) is computed exactly; no stored-scenario
error), and \(L = T + 1 = \tfrac{11}{5}\) --- the kernel difference \(k_t - k_{t'}\) has \(L^1\) norm at most
\(T|t - t'|\) (it is \emph{not} supported on \([t', t]\): the position
rows also differ by \((t - t')n_j\) on \([0, t']\); the integral bound is
what (6) consumes), while \(|\beta_t - \beta_{t'}| = |t - t'|\); the
sandwich (8) thus bounds \(J \le \rho(h) + Th/4 + \tfrac{11}{5}
h_t\), while the witness pair proves the realized gap is exactly
\(Th/4\). The \(\max \bar{e}_a\) form of (8) is load-bearing: with the
loose but verified choice \(\bar{e} = e + \tfrac12\) the program returns
\(\rho = -\tfrac{47}{100}\) at \(h = \tfrac1{10}\), and the certified bound \(J \le \rho + \bar{e} + \delta_\beta + L h_t = \tfrac{17}{100}
\ge \tfrac3{50}\) stays valid, while the uniform special case
\(R_U V_U h_u\) would claim \(J \le -\tfrac{6}{25}\) --- false at
\(J = \tfrac3{50}\). The domain is larger than the campaign: for every uniform aligned mesh
(\(h\) dividing both \(\tau\) and \(T - \tau\)) the same primal--dual
witness pair proves \(\rho(h) = \tfrac{3}{50} - Th/4\) --- the dual
objective \(\sum_k a_k\) telescopes to \((T-\tau)^2/2\) independently of
\(h\) --- confirmed exactly at the sixth aligned mesh \(h = \tfrac1{15}\),
where \(\rho = \tfrac1{25}\) (\(Q = 48\), \(687\) inequalities). No
floor is claimed off the uniform family: on an aligned mesh with unequal
blocks the endpoint moment error is \(\sum_J |J|^2/4\), and the allowed
two-block mesh \([0, \tfrac15], [\tfrac15, \tfrac65]\) gives
\(\rho = \tfrac3{50} - (\tfrac1{100} + \tfrac14) = -\tfrac15\),
certified by the same witness pair (primal row
\(-\tfrac12 + \tfrac{14}{25} - \tfrac{26}{100}\); dual
\(\tfrac12 + (\beta + \tfrac{26}{100}) = \tfrac15 > 0\)). For affine
kernels the block moment error is exactly \(|J|^2/4\) per block, so on
this instance \(e = \sum_J |J|^2/4\) on \emph{every} mesh, uniform or
not, and equal blocks minimize \(\sum_J |J|^2\) at fixed count:
time-refinement here has no conflict structure to target. The
conflict-interval decomposition (12) is a different axis --- it prices
\emph{information} cell splits, not control-block splits --- and it is
the axis the ladder subsection below reads prescriptively. All dimension counts
(\(mQ + 1\) variables; \(Ms(N_t{+}1) + p_U Q\) inequalities), all witness
inequalities, and the closed-form law are machine-verified in exact rational
arithmetic by the verification scripts of Section~\ref{lp-methods}, which
regenerate Table~\ref{lp-tab:mesh} verbatim and re-run the solver campaign.

\subsection{Computation, complexity, and independent verification}\label{lp-computation}

With \(Q\) information--time blocks and \(N_t + 1\) time labels, the program
has \(mQ + 1\) variables and at most \(Ms(N_t{+}1) + p_U Q\) inequalities
(\(p_U\) input-polytope facets); it is polynomial-time in these dimensions and the encoded rational data
size (weakly polynomial bit complexity) for the exact rational benchmark
programs, and for interval-enclosed inputs once the enclosures are fixed ---
conditional on the validated coefficient bounds and certified row
integrals.
Initial-state dependence is affine, so no vertex enumeration of \(P_\theta\)
is required: at each label solve
\(\max_{\xi \in P_\theta} c_i^\top X_\theta(t,0)\xi\) and store a maximizing
or near-maximizing point --- the verifier needs only membership of the stored
point in \(P_\theta\), not global optimality; near-optimality bears on the
convergence guarantee, not on a positive certificate's validity. The same
holds for stored disturbance signals.

For rational outer rows \(\widehat A v \le \widehat b\),
\(F_{\mathrm{blk}} v \le f_{\mathrm{blk}}\), a residual-bounded certificate
test suffices: with \(\lambda, \mu \ge 0\),
\(r = \widehat A^\top \lambda + F_{\mathrm{blk}}^\top \mu\),
\(c = \lambda^\top \widehat b + \mu^\top f_{\mathrm{blk}}\), the condition
\[
c + \sum_j h_U(-r_j) < 0
\tag{11}
\]
is sufficient, since any feasible \(v\) satisfies
\(-\sum_j h_U(-r_j) \le r^\top v \le c\): the upper bound is the
multiplier combination, the lower bound is termwise minimization over
\(U\), and the two displayed inequalities make the sign convention
auditable. The actuator rows \(F_{\mathrm{blk}} v \le f_{\mathrm{blk}}\)
already represent \(v_j \in U\); their multipliers in the residual test
are redundant but legitimate, and the verifier uses both paths as
independent cross-checks. All integration, rounding, and
model-enclosure errors must be included before applying the test, and
uncertain row data must be relaxed toward a larger feasible controller set:
overestimating the worst disturbance contribution can create a false
obstruction --- the sign convention of the \(\beta_a^{+} \ge \beta_a\)
relaxation of Section~\ref{lp-bridge}: a nonmaximizing stored disturbance
gives a larger \(\beta_a^{+}\), hence a weaker constraint, which is the
sound direction for a lower bound on \(J_{\mathcal{I}}(T)\). A deletion-based sparsification attempts supported labels for deletion
one at a time in a fixed order, keeping a deletion only if the prescribed
target margin survives the feasibility re-solve; it terminates with an
irredundant (not claimed minimal) supported label set, at one feasibility
solve per safety row; exact rank compression requires verified rank, and
floating-point singular-value truncation must be charged to the error budget.
If \(J_{\mathcal{I}}(T) = \gamma > 0\), a mesh with
\(\max_a \bar{e}_a + R_D V_D h_d + \delta_{\mathrm{init}} +
\delta_{\mathrm{num}} + L h_t < \gamma\) guarantees detection --- an explicit margin-dependent resolution bound
(mesh resolution as a function of the margin \(\gamma\)), not a
certificate-size bound in general; the uniform choice \(\bar{e}_a
\equiv R_U V_U h_u\) gives the announced special case.

\subsection{Information--time rank: input dimension does not bound the witnesses}\label{lp-rank}

\begin{proposition}[arbitrarily large minimal obstructions with one input]\label{lp-prop:rank}
For every integer \(q \ge 1\) there is a continuous-time system with scalar
input \(u(t) \in [0,1]\) and \(q+1\) indistinguishable modes such that every
proper subbelief is viable, the full \((q{+}1)\)-mode belief is not, and
every obstruction certificate must involve all \(q+1\) modes.
\end{proposition}

\begin{proof}
Let \(T = q\), \(I_j = [j-1, j]\). Mode \(j \le q\) carries the scalar state \(x\) obeying
\(\dot x = \mathbf{1}_{I_j}(t)(u - 1)\), \(x(0) = 0\); mode \(q+1\)
carries \(\dot x = -u\), \(x(0) = q - \varepsilon\), \(0 < \varepsilon
< 1\). Require \(0 \le x \le q\) --- the terminal bounds are the binding
ones, each \(x\) being monotone along its mode --- and there are no
observations. Writing
\(w_j = \int_{I_j} u\), safety requires \(-w_j \le -1\) for \(j \le q\) and
\(\sum_j w_j \le q - \varepsilon\); summing gives \(0 \le -\varepsilon\), a
contradiction. If mode \(q+1\) is removed, \(u \equiv 1\) is safe; if mode
\(j \le q\) is removed, \(u = 0\) on \(I_j\) and \(u = 1\) elsewhere is
safe. A certificate supported on fewer than \(q+1\) modes would also
certify the corresponding proper subbelief, contradicting its viability, so
every obstruction certificate involves all \(q+1\) modes. Equal weights give certificate margin \(\varepsilon/(q+1)\); the
kernel rank is \(q\), and \(q+1\) witnesses are necessary.
\end{proof}

A statement such as ``at most \(m+1\) compatible states suffice to witness
nonviability'' (Helly's theorem, Helly, 1923) is valid for convex common-action sets in \(\mathbb{R}^m\) but
false for a common blind control function: the correct dimension counts
independent temporal and informational control decisions. This distinction
is consequential for review intervals, delayed sensing, and
certificate-size claims; its discrete counterpart is a scoping caveat of the
finite audits of Abaee (2026, Worked systems for the obstruction calculus).

\subsection{Refinement erosion and oracle recourse}\label{lp-erosion}

Suppose a coarse information cell \(C\) is split into finer cells
\(C_1,\dots,C_k\). Because
\(\phi_U(g_1 + \dots + g_k) \ge \sum_j \phi_U(g_j)\), the same witness
satisfies \(\Gamma_{\mathrm{coarse}} \ge \Gamma_{\mathrm{fine}}\), with the
exact difference
\[
\int \Bigl[ \phi_U\Bigl(\sum_j g_{C_j}\Bigr) - \sum_j \phi_U(g_{C_j}) \Bigr]
dt.
\tag{12}
\]
For \(U = [-\bar u, \bar u]^m\) this is
\(\bar u \int \bigl[\sum_j \|g_{C_j}(t)\|_1 - \|\sum_j g_{C_j}(t)\|_1\bigr] dt\):
the certificate identifies the precise time intervals on which
indistinguishable branches demand conflicting control effort. Equation (12)
is the change in the same witness's margin under refinement, not a claim
about the optimized viability value: a coarse cell can strengthen an
obstruction, and splitting it can dissolve the certificate while the
system's verdict is unchanged. The inequality is proved above by pointwise minimization; on rational
instances it is additionally exactly checkable, and the verification record
confirms \(\phi_U(g_1 + g_2) \ge \phi_U(g_1) + \phi_U(g_2)\) on
\(3{,}240\) deterministic rational pairs for the hexagonal \(U\) of
Section~\ref{lp-instance} as a software regression test, not as evidence for
the identity. The discrete counterpart of this axis is the master
monotonicity of the worked-systems companion (Abaee, 2026, Worked systems for the obstruction calculus): its three
axes (observations, policies, memory) enlarge kernels exactly as
refinement here tightens the sandwich, and positivity is non-monotone on
both sides.

The exact theorem assumes exogenous mode observations. A useful sound
extension does not: if finitely many admissible initial-state/disturbance
scenarios are provably indistinguishable until \(\tau\), give the controller
after \(\tau\) separate arbitrary control signals for every stored scenario
--- an oracle relaxation of the true policy class. Scenarios are stored
initial-state/disturbance pairs, indistinguishable through \(\tau\) in
the sense that they generate identical observation histories; after
\(\tau\) the oracle reveals the realized scenario, and separate controls
are granted even where scenarios remain observationally identical. The
soundness direction: every real output-feedback policy induces one common
pre-\(\tau\) control, the oracle class contains every such policy, and
nonviability in the enlarged class transfers down. If
\[
\int_0^\tau \phi_U\Bigl(\sum_\ell \lambda_\ell k_\ell\Bigr) dt
+ \sum_\ell \int_\tau^T \phi_U(\lambda_\ell k_\ell)\, dt
- \sum_\ell \lambda_\ell \beta_\ell > 0,
\tag{13}
\]
with \(\beta_\ell\) computed from the explicitly stored scenario, then no
actual output-feedback policy can be safe. This covers arbitrary sensing
after a known blackout; it is sound but generally incomplete.

A third extension replaces exact scenario storage --- and the exact
revelation of the audited model --- by a certified \emph{belief cell}: a
set \(C \subseteq \Theta \times \mathcal{X}\), carried with a polytope
enclosure, that the observation history certifies contains the true
mode--state pair. --- a set-valued observer with a guaranteed error radius, in the sense
of the bounding approaches to estimation (Milanese et al., 1996).

\begin{proposition}[belief-cell certificates]\label{lp-prop:beliefcells}
Endow each row \(a\) with the cell label
\(\beta_a^{C} = \sup_{(\theta,x) \in C} \beta_a(\theta,x) +
\delta_C\), the stored-scenario label replaced by its worst case over the
cell, with \(\delta_C \ge 0\) the certified enclosure radius of the
position part (the position rows are unit). Then:
\emph{(i)} program (5) under the cell labels remains a linear program
(conic wherever cell support functions enter), and the sandwich (8) and
part (ii)'s soundness hold verbatim with the cell-filtration value in the
middle: every policy adapted to the cell filtration is relaxed, and a
positive value certifies nonviability of the original system.
\emph{(ii)} Splitting \(C\) into \(C_1, \dots, C_k\) erodes a cell
witness by exactly (12), charged cell by cell, and refinement to singletons
recovers the stored-scenario labels (13): observation design is the same
refinement axis as mesh design. \emph{(iii)} Margins are radius-bounded on
both sides: a primal witness whose worst row slack is \(s\) remains a
viability witness for every \(\delta_C \le s\), and a dual certificate
of margin \(\Gamma\) survives every \(\delta_C < \Gamma\) (the
labels' weights sum to one, so radius \(\delta_C\) moves the margin to
\(\Gamma - \delta_C\)).
\end{proposition}

\begin{proof}
(i) a cell-adapted policy must keep every certified pair in \(C\) safe, so
the cell-sup rows relax exactly the cell-adapted class from above; the
sandwich argument uses only the relaxation direction and the enclosure
bound, and the upper side's charges are realized by held controls as before.
(ii) the cell label is a pointwise supremum, so splitting replaces each sup
by the maximum of subcell sups; the resulting gap is the displayed
subadditivity \(\phi_U(\sum_j g_{C_j}) - \sum_j \phi_U(g_{C_j})\) of
(12), and singleton cells restore the stored labels of (13). (iii) shifting
the true state inside a radius-\(\delta_C\) ball translates every affine
row value by at most \(\delta_C\): a witness with slack \(s\) absorbs
it, and the dual combination's value drops by \(\delta_C \sum_j
\lambda_j = \delta_C\).
\end{proof}

On the instance the belief cells are the compatible-mode sets, and the
audited delay hierarchy is exactly the cell-refinement family: the full
cell \(\{1,2,3\}\) obstructs past \(\tfrac7{50}\), the two-mode cells
\(\{1,2\}\) and \(\{2,3\}\) are viable to
\(\tfrac{15 - \sqrt{183}}{6}\) and \(\tfrac{10 - \sqrt{58}}{6}\),
and the singleton cells never obstruct (Proposition~\ref{lp-prop:ladder}).
Noisy state cells price in the same units: at \(\tau = \tfrac15\) the
pair viability witnesses tolerate cell radius up to their worst row slacks
\(\tfrac{31}{1250}\) and \(\tfrac{81}{1250}\), and the triple's
obstruction certificate survives any state-cell radius \(\delta_C <
\Gamma(\tfrac15) = \tfrac3{50}\), the margin moving to
\(\tau - \tfrac7{50} - \delta_C\). On the exact-computation companion's
deadline instance the same law is exact: the radius-\(\delta_C\) cell
inherits the viability verdict precisely for
\(\delta_C \le z_0 - (1 + T/2)\), with the boundary cells zero-slack
(joint check archived; Abaee, 2026, Exact belief-state computation at scale II).

\subsection{Reading the certificate prescriptively: the shared-authority ladder}\label{lp-ladder}

A positive certificate is not only a verdict; its weights, kernels, and
facet multipliers encode what a controller would have to change. On the
instance the decoding is exact end to end.

\begin{proposition}[shared-authority ladder]\label{lp-prop:ladder}
\emph{(i)} The pooling weights satisfy \(\sum_j \lambda_j n_j = 0\)
with \(\lambda_j > 0\); hence for every common control \(w \in U\)
some branch has \(n_j^\top w \ge 0\): no common control brakes all
three branches, and the pooled blind-window kernel --- the certificate's
zero charge on \([0, \tau]\) --- is exactly this unbrakeability.
\emph{(ii)} The optimal shared braking authorities
\(t^{*}(B) = \max_{w \in U} \min_{j \in B}(-n_j^\top w)\) are
\(t^{*}(\{1,2\}) = t^{*}(\{1,3\}) = \tfrac25\), attained at
\(u = (-\tfrac25, -\tfrac45)\) and its \(v_2\)-reflection, and
\(t^{*}(\{2,3\}) = \tfrac35\), attained at \(u = n_1\): the
blind-window controls of the pair policies of
Section~\ref{lp-instance} are exactly these optima. \emph{(iii)} Braking at
authority \(t^{*}\) through the blind window and \(u = -n_j\)
afterward holds the critical-position peak at
\(\tfrac{93}{50} + (1 - t^{*})\tau +
\tfrac{(t^{*})^2 - t^{*}}{2}\tau^2\), increasing in \(\tau\) on
\([0, 1/t^{*}]\); the pair delay thresholds --- where the peak reaches
2 --- are \(\tau_{12} = \tau_{13} = \tfrac{15 - \sqrt{183}}{6}
\approx 0.245\) and \(\tau_{23} = \tfrac{10 - \sqrt{58}}{6} \approx
0.397\), against the triple's \(\tau_{\max} = \tfrac7{50}\), with
singletons unbounded (the full-information peak \(\tfrac{93}{50}\) is
delay-independent). For every \(\tau\) at or below its threshold the
critical position stays \(\le 2\) with the velocity rows strict; the verification script verifies the full seven-facet scan exactly at rational
surrogates below each threshold (\(\tau = \tfrac6{25}\) and
\(\tau = \tfrac{39}{100}\)).
\end{proposition}

\begin{proof}
(i) is the pooling identity itself: \(\sum_j \lambda_j
(n_j^\top w) = \bigl(\sum_j \lambda_j n_j\bigr)^\top w = 0\)
cannot have every summand strictly negative. (ii) The objective
\(\max(n_i^\top w, n_j^\top w)\) is piecewise linear, so its
minimum over the hexagon is attained at a vertex or on the
equal-projection segment \((n_i - n_j)^\top w = 0\). On that segment
both authorities coincide: \(w = -c\,(1,2)\) for \(\{1,2\}\) (the
facet \(|u_2| \le \tfrac45\) binds at \(c = \tfrac25\)) and its
mirror for \(\{1,3\}\); \(w = (c, 0)\) for \(\{2,3\}\) (the
vertex \(n_1\) binds at \(c = 1\)); the vertices themselves give at
most \(\tfrac7{25}\) for the mixed pairs and \(\tfrac35\) for
\(\{2,3\}\). (iii) Along the optimal schedule the window exit state is
\(\tfrac{34}{25} + \tau - t^{*}\tau^2/2\) at velocity
\(1 - t^{*}\tau\) on the critical ray, and full braking adds
\((1 - t^{*}\tau)^2/2\); collecting terms gives the displayed
quadratic, whose \(\tau\)-vertex sits at \(1/t^{*} \ge 1\), so it
increases on the audited range, and the thresholds are its exact roots
in \(\mathbb{Q}(\sqrt{183})\) and \(\mathbb{Q}(\sqrt{58})\)
\(\bigl(6\tau^2 - 30\tau + 7 = 0\) and \(6\tau^2 - 20\tau + 7 =
0\bigr)\). The facet scans are exact rational evaluations at the stated
candidates.
\end{proof}

The prescriptive reading follows directly. The obstruction is removed by
shortening the observation window below the relevant rung (the pairs
already tolerate \(\tau = \tfrac15\) under shared braking alone ---
the audited delay is strictly inside both pair rungs), by adding
observers (each information refinement deletes rungs, and the
conflict-interval decomposition (12) prices exactly which splits pay), or
by reallocating authority where the certificate charges it (the post-window
controls already brake at the hexagon's face \(-n_j \in U\), so the
binding resource is the shared window, not the post-observation actuator).
The dual witness agrees: the \(\eta_j\) place their mass on the braking
facets \(\mp n_j\), and the primal witness brakes every mode to rest ---
the certificate is the blind window's unbrakeability, propagated to the
deadline. The decode is quantitative: each redesign lever carries an exact
rate.

\begin{proposition}[redesign sensitivities]\label{lp-prop:redesign}
On the instance, dilate the post-revelation control set to \(aU\) (the
audited hexagon is \(a = 1\)) and let the observation time \(\tau\)
vary. \emph{(i) Timing:} \(\Gamma_a(\tau) = \tau - \tfrac7{50} -
\tfrac{a-1}{2}\): advancing observation by \(\delta\) buys margin
\(\delta\), at every \(a\). \emph{(ii) Post-revelation authority:}
the exact delay threshold is \(\tau_{\max}(a) = \tfrac7{50} +
\tfrac{a-1}{2}\) (equality viable: the safety set is closed), so each
unit of dilated braking authority buys exactly half a unit of delay
tolerance. \emph{(iii) Shared-window authority:} the pooled
pre-observation kernel is \(\sum_j \lambda_j k_j \equiv 0\) as a
vector identity and the witness's pre-observation stationarity is the
pooling identity alone, so the margin is invariant under \emph{any}
enlargement of the blind-window control set, however large: no shared
pre-observation authority moves the threshold. \emph{(iv) Velocity
budget:} the printed dual places no multipliers on the velocity rows
(their primal slacks are strictly positive at the witness), so any
relaxation of \(|v_1|, |v_2| \le \tfrac65\), up to deleting the rows
entirely, leaves the witness feasible with value \(\tau -
\tfrac7{50}\): the obstruction is positional, carried by braking
liability against position facets.
\end{proposition}

\begin{proof}
Under \(U_a = aU\) the support values scale,
\(\max_{u \in aU}\langle u, n_j \rangle = -a\), so the
post-observation charge in (10) is \(-\tfrac{a}{2}\) and
\(\Gamma_a(\tau) = -\tfrac{a}{2} - (\tfrac{16}{25} - \tau - 1) =
\tau - \tfrac7{50} - \tfrac{a-1}{2}\); dual feasibility is preserved
on the dilated facets \(F_a = aF\), \(f_a = af\) by
\(\eta_j/a \ge 0\), which satisfy \(F_a^\top(\eta_j/a) = -n_j\)
and \(f_a^\top(\eta_j/a) = 1\); the threshold solves
\(\Gamma_a(\tau) = 0\), with equality viable as at \(a = 1\). (iii)
\(\sum_j \lambda_j k_j = (t_* - s)\sum_j \lambda_j n_j = 0\)
independent of the control set, and the printed dual uses no
pre-observation input facets --- stationarity there is
\(\sum_j \lambda_j n_j = 0\) alone. (iv) deleting rows removes no
dual term: the witness's velocity-row multipliers are zero and its value
is unchanged.
\end{proof}

The rates order the redesign options exactly: observation timing at one,
post-revelation authority at one half, shared-window authority and the
velocity budget at zero.

\begin{figure}[t]
\centering
\includegraphics[width=0.94\linewidth]{figs_comp2/fig_hierarchy.pdf}
\caption{The delay hierarchy under authority dilation:
\(\Gamma_a(\tau) = \tau - \tfrac7{50} - \tfrac{a-1}{2}\) crosses zero
at \(\tfrac7{50}\), \(\tfrac{39}{100}\), \(\tfrac{16}{25}\) for
\(a = 1, \tfrac32, 2\) (rate one in \(\tau\), rate one half in
\(a\)); the pair delay thresholds
\(\tfrac{15-\sqrt{183}}{6}\) and \(\tfrac{10-\sqrt{58}}{6}\) mark the
shared-braking rungs of Proposition~\ref{lp-prop:ladder}. Same assertion
script.}
\label{lp-fig:hierarchy}
\end{figure}

\subsection{The reference library}\label{lp-library}

The finite side of the pipeline is packaged as \texttt{viacert}, a
standard-library-only Python package (no third-party dependencies) using
exact integer and rational arithmetic.
Its design follows three principles. \emph{Exactness:} the finite audits and exact certificate verifiers use
integer and rational arithmetic; floating-point solver output is never used
as a proof object.
\emph{Fidelity to the verified recursion:} the finite-horizon selector is the
calculus's backward recursion (Abaee, 2026, An obstruction calculus) with the no-repeat protocol
threaded as a parameter, matching the audited implementations rather than a
reformulation. \emph{Pluggability:} new systems enter through a
\texttt{FiniteSystem} interface (states, actions, safe set, transition, null
action); the audited systems --- the shared two-coordinate audit system, the
continuous benchmark caps, the certainty-equivalence law pair, and the
review grid --- ship as default instances.

The public surface comprises the recursion module (viability verdicts,
kernels, and step-one witness extraction), the monitoring module (safe and
common safe-action sets, the partition-form adequacy converse, maximal
common-action sets, the partition census, the review-timing identity with
its threshold convention, and the certainty-equivalence drift audit), and a
command-line self-test of twelve audited identities --- the four-cell kernel
counts \(24/26/25/28\) with the meet identity, the census \(15\) total and
\(7\) adequate partitions with maximal sets overlapping in \((2,2)\), the
benchmark \(Y^{*} = 27/5\), the drift bounds \(21/100\) and \(0\), the
48-cell delay identity with its boundary cell, and witness extraction ---
each a pointer to the published audits (Abaee, 2026, An obstruction calculus; Abaee, 2026, Worked systems for the obstruction calculus). The
inherited counts are software self-tests pointing to the published audits;
they are not used in this paper's main theorem or example, and their full
context lives in the package documentation. In the certification workflow
the library plays the enumerate-and-check role:
finite instances are enumerated exactly, candidate designs are tested
against the fibre criteria, and every claim is archived with the script that
regenerates it. The library is not a reimplementation of this paper's bridge theorems: it carries the exact finite machinery they reduce from, while the bridge's programs live in the campaign scripts of Section~\ref{lp-methods}; consolidating those scripts as a \texttt{viacert.continuous} module is the library's recorded next step, together with three recorded tooling directions: warm-started incremental refinement (re-solve on block splits seeded with the coarse dual, benchmarked against cold re-solve), an arbitrary-precision dual re-solve fallback for interval-coefficient programs (solve at working precision, detect dual-gap tightness, re-solve the tight dual in extended precision), and conflict-targeted block refinement where kernels are non-affine. Consolidation supersedes implementation paths only in future releases:
the archived scripts and their audited identities remain the exact versions
used for this article's verification record. The title's ``reference
library'' is read as this finite-side package plus the certified campaign
pipeline.

\subsection{Scope delimitations}\label{lp-scope}

The exact theorem does not cover: control-dependent observations or active
sensing; arbitrary continuous-state measurements; nonlinear dynamics without
additional validated enclosures; nonconvex actuator sets without a relaxation gap (convexity is
structural: block averages of admissible controls lie in
\(\mathrm{conv}\,U\), and the weak-* compactness of
Proposition~\ref{lp-prop:value} uses that \(U\) is a compact convex
polytope); rate-constrained actuators without modifying the
approximation; or unrestricted infinite-horizon completeness. For general
sensing after a blackout, the oracle-recourse certificate is sound but
incomplete --- exactness in the audited instance belongs to the complete
revelation at \(\tau\), not to the relaxation. For nonlinear systems, a uniform validated flow-error bound can
be added to the row budgets, yielding a sound extension, not automatic
completeness. Against the alternatives: full-information viability is necessary but not sufficient for the
information-constrained question; initial-admissibility and tangency tests
are silent on the instance's obstruction, which is temporal-informational
rather than pointwise-geometric; safety-before-first-observation tests are
silent when a common safe blind control exists; exact finite safety games
exist as discrete models, and what is needed to use them as sound
abstractions of the present continuous systems is a justified
continuous-to-finite bridge, which is what Theorem~\ref{lp-thm:bridge}
provides; and na\"ive held-control infeasibility
excludes only the held-control class. No general-purpose
computational-calculus claim is made: the theorem is a certificate generator
for the stated class, with completeness only under refinement and
semidecidability only with effectively computable bounds.

\subsection{
Conclusion}\label{lp-conclusion}

\emph{The general finding first, because it is the one that survives the instance.} The cost of
certifying a system under incomplete observation is not governed by the size of the system, and it
is worth saying precisely what it \emph{is} governed by, because the two are routinely conflated.
The certificate produced by the pipeline of this paper has a size that scales with the
\emph{information--time product}: the product of how much the controller can learn and over how
long a horizon it must act. It does not scale with the input dimension --- the dimension of the
state, of the disturbance, or of the discretisation standing in for either. A high-dimensional
system about which little can be learned over a short horizon is cheap to certify; a
low-dimensional system observed richly over a long horizon is expensive. Dimensionality is the
wrong axis, and pipelines organised around reducing it are optimising the wrong quantity. This is
the paper's transferable claim, and it is a claim about the \emph{shape} of the problem rather
than about any instance of it. The same conclusion is reached independently, in a different
language and for a different certificate form, in the measure-dual companion (Abaee, 2026,
minimax dual certificates), where the witness size is set by the control dimension \(k\) with a
tight \(k+1\) bound. Two different constructions, one moral: \emph{obstruction certificates are
small objects whose size is fixed by the decision and information structure, not by the size of
the world.}

\emph{What that claim rests on.} A sound and independently checkable certification pipeline under
incomplete observation rests on three requirements, and the bridge theorem delivers all three.
First, the finite controller class must be an \emph{outer} relaxation while the stored adversarial
tests correspond to genuinely admissible scenarios --- relax the controller, not the adversary, or
the certificate certifies nothing. Second, optimality must be certifiable independently of the
solver: the solver's output is a confirmation, not a proof object, and the exact witnesses are
what close the gap. Third, and this is the requirement above, certificate size must be understood
in the information--time product rather than the input dimension.

\emph{What the instance shows.} The three-branch instance shows the resulting certificates
detecting failure that full-information viability, blind-window exit tests, and held-control
infeasibility all miss, and the campaign shows the solver returning the independently proved
optima to numerical precision while exact witnesses close every gap. The reference library carries
the exact finite machinery to new systems.

\emph{What remains open.} Completeness beyond refinement:
 an inner
approximation of the policy class that is computationally effective for
general sensing is not known, and the oracle-recourse extension is the
current sound surrogate.

Three further extension directions are recorded, each a construction of
its own rather than an increment of the present framework; all lie outside
the audited class delimited in Section~\ref{lp-scope}. The cell-conditioned
lift of labels, kernels, and costs for general measurement channels --- the
passage from history partitions to general measurement maps --- is carried
in Proposition~\ref{lp-prop:beliefcells}, with the instance's compatible-mode
hierarchy and radius bounds as the audited instantiation.
\emph{(i)} Stochastic disturbances: replacing support-function rows with
second-moment bounds turns the finite program into a conic one
(second-order cone or semidefinite), with exact rational witnesses
replaced by conic dual certificates; the exact rational optima of the
present campaign remain the benchmark such relaxations must meet.
\emph{(ii)} Decentralized and networked controller sets, infinite-horizon
and periodic regimes, and nonconvex or piecewise-smooth actuators: each
alters a structural input of the bridge theorem (the weak-* compactness
of label strategies, finite-horizon completeness, or the convexity of
the actuator set), and is therefore a new construction rather than a
parameter change.
\emph{(iii)} A cross-tool comparison study that treats the audited
instance as a common benchmark --- exact finite safety games,
interval-reachability toolchains, SMT-backed barrier synthesis --- with
the certificate file as the exchange format; the tooling directions
recorded in Section~\ref{lp-library} and the independent
certificate-file parser of Section~\ref{lp-methods} are its computational
prerequisites.

\subsection{Verification methods}\label{lp-methods}

Every reported mathematical benchmark value and every table entry is
regenerated by deterministic scripts, exact rational arithmetic unless
stated as solver output, standard library plus the solver interface only.

\noindent\textbf{A mechanized layer, distinct from the scripts.} The
checks catalogued below are computations; they do not prove. Separately,
the Farkas core on which the exact primal--dual witnesses rest is
\emph{proved} in a mechanized layer (Lean~4, 60 modules under \texttt{Formalizations/}, no axioms and no admitted gaps): the soundness direction --- the one the certificates actually use --- is derived from the ordered-field interface alone.
\emph{The project is pinned to Lean~4 in \texttt{lean-toolchain}; the checks were run under
\texttt{v4.14.0}, and on 2026-09-30 the source-level invariants were re-verified against the
current tree --- a comment-stripped scan of all \texttt{.lean} files finds no \texttt{sorry},
\texttt{admit}, \texttt{axiom} or \texttt{constant} declaration anywhere. The build has not
been re-run since the toolchain pin moved.} The converse direction is not proved there, and is not
used: its classical argument rests on a separation theorem the layer
does not develop. Three results of this paper remain out of reach of
that layer by construction, and are cited rather than formalized: the
value identity, the bridge theorem, and the belief-cell proposition,
each of which calls for analysis (exponentials and logarithms, or
weak-\(*\) compactness) that an ordered-field interface does not
supply. That is a statement about the formalization, not about the
results.
Three scripts and one package cover the claims, with one master entry
point chaining them. The exact
verification record re-derives the instance's certificates --- delay
threshold, margins, witnesses, erosion identity, rank construction ---
in 51 exact checks. The campaign script rebuilds the program (5) at each of
the five meshes, runs HiGHS, and certifies optimality by the exact
primal--dual witnesses, closing the duality gap; it re-derives the mesh law
and all dimension identities (nine checks). The library self-test runs the
twelve audited identities. A master script chains the three, then
recomputes independently: the mesh values from the closed law
\(3/50 - Th/4\) (uniform aligned meshes), the Farkas contradiction
\(-\tfrac{53}{100} + \tfrac12 = -\tfrac{3}{100}\), the primal row value
\(\tfrac{3}{100}\), the margins \(\tfrac{3}{50}\) and \(\tfrac{3}{100}\),
the moment error \(e = \tfrac{3}{100}\), the relaxed right-hand side
\(-\tfrac{53}{100}\), the delay threshold \(\tau^{*} = \tfrac{7}{50}\), the blind-window safety value \(\tfrac{39}{25}\), and every headline
number in the text; independence here means a separate formula path from
the subordinate scripts, not a separate implementation --- the independent
certificate-file parser is shipped (\texttt{certificate\_exchange\_v1.py})
and validates the archived witnesses of this paper and of the
dual-certificate companion (Abaee, 2026, Minimax dual certificates) exactly, in either direction,
alongside the tooling directions of Section~\ref{lp-library}. The two
figures are generated by \texttt{paper2\_comp\_certification\_figures\_v1.py}
(recompute-then-assert: every plotted peak, threshold, and crossing is
recomputed in exact rational arithmetic and asserted against the printed
values before the figure is written; ten assertions). Timings are single runs on the verification
machine (Linux x86\_64, Intel Xeon @ 2.60 GHz, Python 3.13, SciPy 1.17
HiGHS dual-simplex interface, solver defaults otherwise).

\section{Related work}
\label{lp-priorart}

This paper builds a finite, solver-attainable, independently checkable certificate for a
continuous-time safety value under partial observation. Each of those four properties has an
established literature behind it, and this section says where the present construction sits
relative to each.

\subsection{Viability kernel computation, and the sandwich}

The computation of viability kernels by discrete approximation is the classical route.
Saint-Pierre's viability kernel algorithm (Saint-Pierre, 1994) approximates the finite-horizon
kernel by a backward recursion on a grid, with a convergence statement: for open constraint
sets, the union of the approximating sets over vanishing time step recovers the kernel.
Level-set formulations compute the same objects through Hamilton--Jacobi equations
(Mitchell, Bayen and Tomlin, 2005). Because these Eulerian schemes grid the state space,
their cost grows exponentially with dimension, which motivated Lagrangian reformulations in
terms of reachable sets that scale considerably further (Maidens, Kaynama, Mitchell, Oishi
and Dumont, 2013).

\textbf{The two-sided structure is not new, and this paper does not claim it is.} The
reachability-based viability computations produce a guaranteed \emph{under}-approximation of
the kernel together with a free \emph{over}-approximation, and use the latter to bound the
error of the former (Maidens et al., 2013, Section~3). The present paper's sandwich,
\[
\rho_{\mathcal{H},G}\;\le\;J_{\mathcal{I}}(T)\;\le\;\rho_{\mathcal{H},G}+\bar{e}+\delta_{\beta}+L h_{t},
\]
is of that kind: a computable quantity, and a certified inflation that bounds the true value
from the other side. What differs is the object being sandwiched and the approximations that
generate it. The classical schemes discretise a \emph{state space} and approximate a
\emph{kernel}, a set of states. Here the discretisation is over an outer moment relaxation of
the admissible \emph{information-adapted controls} together with an inner adversarial
selection of finitely many stored labels and scenarios, and the quantity bounded is a
safety \emph{value} over an information state — not a set of initial conditions. The
partial-observation structure, and the fact that the certificate must hold against every
measurable policy rather than against one synthesised controller, is what forces that
change.

\subsection{Moment relaxations}

Relaxing a hard optimisation problem into a sequence of convex problems whose optimal values
converge monotonically to the true optimum is the moment--SOS programme (Lasserre, 2001;
Parrilo, 2003), whose hierarchy produces nondecreasing lower bounds with asymptotic
convergence and, importantly, \emph{checkable} conditions — rank and flat-extension criteria
on the moment matrix — under which a relaxation is certified exact and minimisers can be
extracted. That certification idea is the direct ancestor of the present one.

The differences are concrete. The moment--SOS hierarchy applies to polynomial optimisation
over semi-algebraic sets and is solved by \emph{semidefinite} programming; the relaxations
here are \emph{linear} programs over moment variables for a linear system, and the programs
are solved by a standard LP solver. Convergence in the hierarchy is in the relaxation
degree; convergence here is along a refinement sequence whose mesh and stored-scenario and
enclosure widths vanish (Theorem~\ref{lp-thm:bridge}). And the certified object differs: an
exactness certificate in the moment--SOS setting licenses the extraction of a global
minimiser, whereas here a strictly positive certified margin is an \emph{obstruction}
certificate — a proof of nonviability, not the recovery of an optimum.

\subsection{The scenario approach, and why the guarantee here is not probabilistic}

Sampling a finite set of uncertainty instances and solving the resulting finite program is
the scenario approach (Calafiore and Campi, 2005, 2006), with sharp sample-complexity
results (Campi and Garatti, 2008; Campi and Garatti, 2011, on sampling-and-discarding; Campi,
Garatti and Prandini, 2009, for the systems-and-control setting). The inner approximation
here — finitely many stored admissible labels and scenarios — is superficially the same
move.

It is not the same guarantee. The scenario approach is \emph{probabilistic}: for $N$ sampled
constraints the returned solution violates the unsampled constraints with probability at most
$\epsilon$, with confidence $1-\beta$, where $N$ is chosen to deliver that pair. The
guarantee here is \emph{worst-case}: the scenarios are selected adversarially rather than
sampled, and the bound holds for every measurable information-adapted policy, with no
measure on the uncertainty and no confidence parameter. A reader who imports intuitions from
the scenario literature should note that $\epsilon$ and $\beta$ have no counterpart here;
the error terms $\bar{e}$, $\delta_{\beta}$ and $L h_t$ are deterministic inflation
certificates, not violation probabilities. (The notational collision on $\beta$ is
unfortunate and is flagged rather than resolved.)

\subsection{Verified numerics, and the role of the solver}

Enclosing the solution of a continuous problem in rigorously bounded intervals is the
programme of interval arithmetic and validated ODE integration. This paper deliberately does
not use it. Every exact verification path for the reported rational benchmark runs in integer
and rational arithmetic, and optimality is certified by independent primal and dual witnesses
checked exactly. \textbf{The solver is not a proof object}: the numerical optimum returned by
the LP solver is a confirmation that agrees with the exact rational value to within
$4.2\times10^{-17}$ on the reported cases, and the certificate rests on the exact witnesses,
not on the floating-point output. This is a statement about the division of labour, not a
claim that interval methods are inferior; they address problems this pipeline does not.

\subsection{What is claimed}

Given the above, the contribution is five things.

\begin{enumerate}
\item \textbf{A continuous-to-finite bridge for information states}: an outer moment
  relaxation and an inner adversarial discretisation combine into a finite LP whose value and
  whose error inflation bound the continuous-time safety value on both sides, with
  convergence along any admissible refinement sequence (Theorem~\ref{lp-thm:bridge}).
\item \textbf{Dual-feasibility as a universal certificate}: every dual-feasible solution of
  the finite program yields a finite lower-bound certificate covering \emph{every} measurable
  policy, and a strictly positive certified margin is a finite obstruction certificate.
\item \textbf{The witness count is governed by information--time rank, not input dimension}:
  whenever a positive certificate exists one can be chosen with at most $r+1$ safety labels
  at rank $r$, and a matching construction with scalar input produces obstructions whose
  minimal certificates require arbitrarily many witnesses, saturating the bound
  (Proposition~\ref{lp-prop:rank}). The negative half — that a single input does not bound the
  witness count — is the part most likely to surprise.
\item \textbf{The certificate is prescriptive}, decoding into an exact shared-authority
  hierarchy and into redesign sensitivities that separate what buys tolerance from what buys
  nothing (Propositions~\ref{lp-prop:ladder} and~\ref{lp-prop:redesign}).
\item \textbf{Exact verification paths by construction}: primal and dual witnesses checked in
  rational arithmetic, with the solver demoted to a numerical confirmation.
\end{enumerate}

\noindent Not claimed: the sandwich idea as such (Maidens et al. and the reachability
literature), moment relaxation as such (Lasserre; Parrilo), scenario discretisation as such
(Calafiore and Campi; Campi and Garatti), or interval enclosure. The novelty, such as it is,
lies in assembling these into a certificate that is valid against every measurable
information-adapted policy under partial observation, and in the rank result of item~3.

\part*{Part II --- The measure dual: what the obstruction is dual to}
\label{part:measure-dual}
\addcontentsline{toc}{part}{Part II --- The measure dual}

\maketitle

\begin{abstract}
The common-action obstruction certifies that no single instrument serves
every branch of an information set; its polyhedral form carries a Farkas
multiplier (Abaee, 2026, An obstruction calculus, Section 3.2). This paper states and proves the
measure dual of the same obstruction: for control-affine dynamics with a
compact convex control set, the common safe-action set is empty precisely
when an adversarial probability measure on the active
boundary--disturbance bundle drives the expected inward drift strictly
negative against every control, with a certificate supported on at most
$k+1$ points ($k$ the control dimension; the bound is tight). The
polyhedral case recovers the $l_1$-normalized Farkas pair exactly; the
singleton case recovers the Isaacs drift with the adversarial
\emph{distribution} over disturbances load-bearing. A two-action instance
exhibits a strict minimax gap: without convexity the equivalence fails and
no measure certifies --- the audited finite codex systems sit outside this
scope by design, certified combinatorially. An obstruction taxonomy
(physical, epistemic, institutional) is instantiated on the audited grid,
and a calibrated two-stock resource model --- one fixed control set
throughout --- yields a closed-form critical aggregate
$Y^{*} = \tfrac{27}{5}$: aggregate-yield fibres are viable exactly up to
$Y^{*}$, the obstruction at $Y = 6$ is certified by a two-point dual
measure with margin $\tfrac{3}{50}$, and the fibre-width formula
$\Delta x_{1}(Y) = Y - 4$ gives the aggregation rule its exact reading.
Dynamic-envelope, quantization, and stochastic-bridge extensions are
recorded as conjectures; the residual dynamic-observation completeness gap
of the source paper remains open. All claims are machine-verified in exact
rational arithmetic (8 check families).
\end{abstract}

\vspace{0.4em}
\noindent\textbf{Keywords:} minimax duality; Farkas certificates; zero-sum
games; obstruction certificates; adversarial priors; Isaacs condition;
exact arithmetic

\section{Introduction}

The obstruction calculus's certificates are emptiness statements: a belief is
obstructed when a correspondence of jointly acceptable responses is empty
(Abaee, 2026, An obstruction calculus, Theorem 2; Abaee, 2026, Worked systems for the obstruction calculus). For polyhedral safe-action sets the
emptiness carries an explicit dual object --- the Farkas multiplier vector (Farkas, 1902)
--- so that the obstruction itself becomes a checkable artifact; the
sparse-witness bound (at most $m+1$ compatible states) and the singleton
Isaacs drift (Isaacs, 1965) complete the static picture. This paper unifies the static layer: the measure dual of the common-action obstruction, of which the normalized Farkas pair, the sparse witness, and the Isaacs reduction are special cases. The measure dual is verified exactly; the sparsity bound is proved by a finite-reduction route (finite reduction, linear-programming duality, basic-solution support); the singleton Isaacs reading is stated with its load-bearing distributional caveat; the equality's convexity boundary is exhibited by counterexample; the state-versus-epistemic separation is adopted as a working axiom; and a calibrated two-stock benchmark with a single fixed control set is solved in closed form. The dynamic layer --- an adversarial-prior envelope for the restricted-information game --- is then formulated and analyzed, and the residual dynamic-observation completeness gap (Abaee, 2026, An obstruction calculus, Open Problem 1) is \emph{not} closed here. The contributions are: the measure dual of the common-action obstruction, exactly verified, with the normalized Farkas pair, the sparse witness, and the Isaacs reduction as special cases, and the equality's convexity boundary exhibited by counterexample (Section 2); the calibrated two-stock benchmark solved in closed form (Section 4); and the dynamic layer --- the adversarial-prior envelope for the restricted-information game, its three obstructions, and the review-sequence envelope's exact couplings (Sections 5--6).

\begin{table}[h]
\centering\footnotesize
\caption{The theorem and its check families: what the verification record
certifies.}
\label{dual-tab:checks}
\begin{tabular}{@{}l p{4.6cm} l@{}}
\toprule
Family & Statement certified & Checks \\
\midrule
G1 & general tower on all trees ($K{=}3$, adaptive filtration); & \\
   & martingale property at every cell, all 8192 policies & \\
G2 & general DP equals path enumeration ($K{=}2$, 19683 policies); & \\
   & $Q_0 = N_1$ equals the tree supremum on both instances & \\
G3 & refinement monotonicity, strict on a crafted functional & \\
G4 & bridge direction: $Q_0 < 0$ defeats every policy measurably & \\
G5 & coupling selection is real: $0$ vs $+1$ at equal marginals & \\
G6 & text, retired markers, proof pairing & \\
\bottomrule
\end{tabular}
\end{table}

\section{The measure dual of the common-action obstruction}

Let the safe set be $V = \{x: q_{j}(x) \ge 0,\ j = 1,\dots,m\}$ with
$q_{j}$ of class $C^{1}$; let $D$ be a compact-valued upper
semicontinuous disturbance correspondence; let the control set
$U \subseteq \mathbb{R}^{k}$ be compact and \emph{convex}; and let the
dynamics be affine in the control,
$f(x,u,d) = f_{0}(x,d) + G(x)u$, with $f_{0}$ and $G$ continuous. For a
compact information set $B \subseteq V$ (a fibre or estimation tube) with
nonempty active bundle
\[
\mathcal{K}(B) = \{(x,j,d):\ x \in B,\ q_{j}(x) = 0,\ d \in D(x)\},
\]
define the inward drift
$\psi(x,j,d;u) = \nabla q_{j}(x)\cdot[f_{0}(x,d) + G(x)u]$ and the
expected drift $\Phi(u,\mu) = \int_{\mathcal{K}(B} \psi \,d\mu$ over
Borel probability measures $\mu$ on $\mathcal{K}(B)$.

\begin{theorem}[minimax dual certificate]\label{dual-thm:dual}
Under the standing hypotheses the common safe-action set
$\mathcal{R}^{B}_{V}(B) = \bigcap_{x \in B}\mathcal{R}_{V}(x)$ is empty if
and only if there exist an adversarial measure
$\mu^{*} \in \mathscr{P}(\mathcal{K}(B))$ and a margin $\varepsilon > 0$
with
\begin{equation*}
\max_{u \in U}\,\Phi(u,\mu^{*}) \le -\varepsilon < 0.
\tag{$\star$}
\end{equation*}
Moreover
$\max_{u \in U}\min_{\mathcal{K}(B)}\psi
 = \min_{\mu}\max_{u \in U}\Phi(u,\mu)$.
\end{theorem}
\begin{proof}
$\mathscr{P}(\mathcal{K}(B))$ is weak-$*$ compact and convex (Banach --
Alaoglu), and $\Phi$ is continuous, affine in $u$ (control affinity), and
linear in $\mu$; Sion's minimax theorem (Sion, 1958) applies, and the inner infimum
over measures of a linear functional is attained by the Dirac measures,
so
$\max_{u}\min_{\mu}\Phi = \max_{u}\min_{\mathcal{K}(B)}\psi
 = \min_{\mu}\max_{u}\Phi$.
For the equivalence: $u \in \mathcal{R}^{B}_{V}(B)$ exactly when
$\min_{\mathcal{K}(B)}\psi(x,j,d;u) \ge 0$; the minimum is continuous in
$u$, so nonemptiness of the common safe-action set is equivalent to
$\max_{u \in U}\min_{\mathcal{K}(B)}\psi \ge 0$ (the maximum is attained
on the compact convex $U$), and emptiness is equivalent to the common
value being strictly negative --- which, by the equality, is equivalent to
some $\mu^{*}$ with $\max_{u}\Phi(u,\mu^{*}) < 0$; the margin is the
slack.
\end{proof}

\begin{proposition}[atomic support; $k+1$ witnesses, tight]
\label{dual-prop:sparse}
If $\mathcal{R}^{B}_{V}(B) = \varnothing$, the certificate $\mu^{*}$ may
be chosen atomic with support of at most $k+1$ triples, and the bound is
tight in dimension $k$.
\end{proposition}
\begin{proof}
Write $\psi(x,j,d;u) = v(x,j,d)\cdot u + c(x,j,d)$ with
$v = G(x)^{\top}\nabla q_{j}(x)$ and $c = \nabla q_{j}(x)\cdot f_{0}(x,d)$.
On the compact set $U \times \mathcal{K}(B)$ the function $-\psi$ is
continuous with
$\min_{u}\max_{\mathcal{K}(B)}(-\psi) = 2\eta > 0$ (Theorem~\ref{dual-thm:dual});
a standard finite-cover argument extracts finitely many triples
$(x_{i},j_{i},d_{i})$, $i \le N$, whose drift rows still satisfy
$\min_{u \in U}\max_{i}(-\psi_{i}) \ge \eta$. Min--max duality for the
finite system (linear programming duality over the simplex) yields
$\lambda \ge 0$, $\sum_{i}\lambda_{i} = 1$, with
$\sum_{i}\lambda_{i}\psi_{i}(u) < 0$ on $U$; an extreme such
$\lambda$ is a basic solution of $k+1$ active constraints and is supported
on at most $k+1$ indices (standard theory: Chv\'{a}tal, 1983, chapters 2
and 7; the pairing argument itself is Carath\'eodory's theorem in
$\mathbb{R}^{k}$). The tightness bound is Helly's theorem in the plane
(Helly, 1923) applied to the row half-spaces. Tightness: in dimension $k = 2$ the rows
$\{u_{1} \ge 1\}$, $\{u_{2} \ge 1\}$, $\{u_{1}+u_{2} \le 1\}$ on the
compact square $[-10,10]^{2}$ are pairwise intersecting with empty total
intersection (checked on the rational grid by the record), so no two
witnesses suffice; for $k = 1$ the instance of E1 certifies with exactly
two atoms.
\end{proof}

\begin{proposition}[recoveries]\label{dual-prop:recover}
(i) \emph{Polyhedral.} If the drift rows are affine with rational data,
the measure certificate is the $\ell_{1}$-normalized Farkas pair of
Abaee, 2026, An obstruction calculus, Section 3.2: on the two-floor instance
$\{u \le \tfrac25\} \cap \{u \ge \tfrac35\}$ with $U = [0,1]$, the
measure $\lambda = (\tfrac12,\tfrac12)$ makes the expected drift
\emph{constant} $-\tfrac1{10}$, the normalized multiplier pair of the
worked certificate in Section 3.4 there; on the feasible contrast
$\{u \le \tfrac25\} \cap \{u \ge \tfrac15\}$ both sides of
Theorem~\ref{dual-thm:dual} equal $+\tfrac1{10}$ and no certificate exists.
(ii) \emph{Singleton.} For $B = \{x\}$ the dual measures live on the
disturbance set; on the exact convex instance
$\psi(u,d_{1}) = 1-2u$, $\psi(u,d_{2}) = 2u-1$,
$\max_{u}\min_{d}\psi = \min_{\mu}\max_{u}\Phi = 0$, while the
single-worst-disturbance reading $\min_{d}\max_{u}\psi = 1$ is
\emph{not} the dual: the adversarial distribution over disturbances is
load-bearing, and the Isaacs drift is its Dirac-supported inner
infimum.
\end{proposition}

\begin{proposition}[convexity boundary]\label{dual-prop:gap}
The equivalence of Theorem~\ref{dual-thm:dual} fails without control
convexity. With $U = \{0,1\}$ and drift rows
$(-1,+1)$, $(+1,-1)$, the common safe-action set is empty, yet
$\max_{u}\min\psi = -1$ while $\min_{\mu}\max_{u}\Phi = 0$: a strict
minimax gap, and no adversarial measure certifies. The audited finite
action codices of the catalogue notes are non-convex by construction;
there the obstructions are certified combinatorially, and the present
theorem is a companion for the continuous, control-affine layer.
\end{proposition}

\medskip
\noindent\textbf{Shared symbols across the companion papers.} The
companion papers share this vocabulary with paper-local meanings; the
letters below carry meanings here that differ there (including this
paper's own three uses of $K$):
\begin{center}\small
\begin{tabular}{@{}p{0.42\linewidth}p{0.50\linewidth}@{}}
\toprule
\textbf{Here} & \textbf{In the companions} \\
\midrule
$K$: the window/tree depth; $\mathcal{K}(B)$: the contact set; $K_{i}$: the benchmark's carrying capacities & the safe set and kernel $K$ (obstruction calculus); the full-information kernel $K^{T}_{\mathrm{full}}$ (certification); the carrying capacity (forecast ladder) \\
$\mu$, $\mu^{*}$: the adversarial probability measures & Farkas multipliers (obstruction calculus); the certification companion's dual measures \\
$\lambda$: the prior over stored branches; the dual weights & the Farkas multiplier and observer decay rate (obstruction calculus); the certificate weights (certification) \\
$\Gamma_{h}$: the recourse certificate & the certificate margin $\Gamma_{\mathcal{I}}$ (certification); the alpha-vector witness sets $\Gamma^{\Pi}_{k}$ (sufficiency) \\
$\rho$: the interpolation parameter of the DR family & the contamination radius $\rho$ (sufficiency); the mesh relaxation parameter (certification); the net biomass multiplier (applied regime viability) \\
$\beta$: the margin & the stored per-label bound $\beta_{a}$ (certification) \\
$\tau$: the reveal time & the observation time $\tau$ (certification); the delay times (obstruction calculus); the worst-case exit time $\tau^{*}$ (worked systems) \\
$\sigma$: the kernel time variable & the floored exit count $\sigma^{*}$ (worked systems); a noise scale (forecast ladder) \\
\bottomrule
\end{tabular}
\end{center}


\section{The obstruction taxonomy}

The standing axiom adopted from the audit: \emph{a dead state is never an
information failure.} Three classes of nonviability are distinguished and,
on the audited grid, are instantiated side by side (all exact):
\emph{physical} --- the corner $(1,1)$ under two floors is nonviable as a
singleton under full information (no instrument keeps both floors);
\emph{epistemic} --- the audited belief $\{(1,2),(2,1)\}$ has every branch
viable and the joint set empty (failure by merge); \emph{institutional}
--- under the unanimity protocol with agency 2 restricted to vote 1, the
decentralized kernel collapses from $12$ to $0$ (failure by protocol;
full enumeration). The taxonomy disciplines attribution across the
catalogue: mode-blind losses on the same grid are generated entirely at the singleton level, and no count
that mixes the classes is admissible as an epistemic claim.

\section{A calibrated two-stock benchmark}

Leaving the audited grid, the continuous layer needs a worked system with
one fixed control set throughout; the two-stock system below is the
running instance of the continuous layer --- the caps-fibre benchmark
shared with the worked-systems companion. Two renewable stocks
($r_{1} = 1$, $r_{2} = \tfrac45$, $K_{1} = K_{2} = 10$, cross-competition
$\tfrac1{20}$, shocks $|d_{i}| \le \tfrac1{10}$), floors
$x_{i} \ge 2$, and the single control set
$U = \{u \ge 0:\ u_{1} + u_{2} \ge 2\}$ (a demand floor; no capacity
constraint is imposed). On the aggregate-yield fibre
$Y = x_{1} + x_{2}$ within the floors, the active-floor harvest caps are
\[
\mathrm{cap}_{1}(Y) = \tfrac32 - \tfrac{Y-2}{10},\qquad
\mathrm{cap}_{2}(Y) = \tfrac{59}{50} - \tfrac{Y-2}{10},
\]
so the fibre is viable exactly when
$\mathrm{cap}_{1}(Y) + \mathrm{cap}_{2}(Y) \ge 2$, i.e.
\begin{theorem}[critical aggregate]\label{dual-thm:benchmark}
$Y^{*} = \tfrac{27}{5}$: aggregate fibres with $Y \le \tfrac{27}{5}$ are
viable and those with $Y > \tfrac{27}{5}$ are obstructed, the value being
exact ($\mathrm{cap}_{1} + \mathrm{cap}_{2} = 2$ at $Y^{*}$). At
$Y = 5$ the fibre is viable with witness $u = (\tfrac65,\tfrac45)$; at
$Y = 6$ it is obstructed (cap sum $\tfrac{47}{25} < 2$) with dual
certificate $\mu^{*} = (\tfrac12,\tfrac12)$ on the two active-floor
atoms --- two points, within the $k+1 = 3$ bound --- and margin
$\varepsilon = \tfrac{3}{50}$. The fibre criterion itself --- viability
exactly on the cap-sum inequality, audited per aggregate with radius
bounds --- is carried in the worked-systems companion (Abaee, 2026, Worked systems for the obstruction calculus);
this section contributes the two-stock derivation and the fibre-width
formula.
\end{theorem}
The reading is the opposite of the intuition that abundance certifies
safety: cross-competition tightens the harvest caps as the aggregate
grows, so richer fibres obstruct first --- aggregate yield is a failing
indicator in both directions. The fibre-width formula makes the
aggregation rule exact: the $Y$-fibre within the floors is
$x_{1} \in [2,\,Y-2]$, of width $\Delta x_{1}(Y) = Y - 4$, so a linear
index cannot separate states closer together than the floor buffer ---
the composite index is certifying exactly when the width stays inside the
buffer, and the certainly-safe reporting of the applications companion
(Abaee, 2026, Robust viability of the 2J3KL limit reference point) governs otherwise.

\section{The dynamic envelope}

A natural candidate formulation of a dynamic certificate suggests itself: an adversarial measure curve $t \mapsto \mu^{*}_{t}$ whose expected-margin envelope $Q(t)$ decreases at a certified rate, governed by an information-set operator acting on the state. The candidate fails, for three reasons that shape the formulation developed below; this section records the reasons, formulates the envelope, and proves what exact methods reach. The residue remains conjectural.

\subsection{Three obstructions to the candidate}

\emph{(i) A Bellman object posed as a pointwise Hamiltonian.} The candidate operator maps a
value along trajectories to an expectation over an information set of
measures --- a Bellman-type, hence nonlocal, object --- while dressed as a
pointwise differential operator. Restricted-information games do not
produce pointwise Hamiltonians in the state; their well-posed home is the
information state (here, the belief), on which the object is a dynamic
program. The nonlocality is not itself the obstruction; the pointwise differential form is.

\emph{(ii) An inconsistent measure curve.} A curve $t \mapsto
\mu^{*}_{t}$ of measures carries no requirement that $\mu^{*}_{t}$ be the
marginal at time $t$ of one measure on disturbance paths. Without that
consistency the envelope has no time consistency: nothing connects $Q(t)$
to $Q(t')$ for $t' > t$, so no backward recursion and no comparison can
even be posed. The formulation developed below replaces the curve by a single adversarial prior $\mu^{*}$ on disturbance paths; conditional expectation then supplies
the martingale structure without any PDE.

\emph{(iii) Expectation versus trajectory.} For the certificate direction
the selection argument is not needed at all: an expectation of a bounded
functional that is strictly negative exhibits a positive-measure set of
realizations on which the functional is strictly negative. Selection
theorems are required only for the converse value-attainment direction,
which the certificate does not use.

\subsection{The envelope formulation}

\begin{lemma}[expectation-to-realization bridge]\label{dual-lem:bridge}
Let $\mu^{*}$ be a probability measure on disturbance paths compatible with
the information structure, let $\pi$ be nonanticipative, and let the
safety functional $F^{\pi}(\omega)$ --- the realized weighted facet margin
at the horizon --- be bounded below and $\mathcal{F}_{T}$-measurable. If
$\mathbb{E}^{\mu^{*}}[F^{\pi}] < 0$, then
$\mu^{*}(\{\omega: F^{\pi}(\omega) < 0\}) > 0$: some compatible disturbance
realization violates the margin, and the set of such realizations carries
positive $\mu^{*}$-measure.
\end{lemma}
\begin{proof}
Nonanticipativity of $\pi$ makes $F^{\pi}$ an $\mathcal{F}_{T}$-measurable
functional of the disturbance path, and boundedness below makes its
expectation well-defined. If $\mu^{*}(F^{\pi} < 0) = 0$, then
$F^{\pi} \ge 0$ holds $\mu^{*}$-almost surely, so
$\mathbb{E}^{\mu^{*}}[F^{\pi}] \ge 0$ --- contradicting the hypothesis. The
realizations of the positive-measure set are compatible by construction,
since $\mu^{*}$ is supported on the compatible set.
\end{proof}

\begin{definition}[adversarial-prior envelope]\label{dual-def:envelope}
Fix a horizon $T$, a bounded safety functional $F$ (positive values safe),
and one consistent adversarial prior $\mu^{*}$ on disturbance paths. The
envelope is
\[
Q(t) \;=\; \operatorname*{ess\,sup}_{\pi} \;
\mathbb{E}^{\mu^{*}}\bigl[\, F^{\pi} \;\big|\; \mathcal{F}_{t} \,\bigr],
\]
the essential supremum over nonanticipative policies of the conditional
expected safety functional. For each fixed $\pi$ the process
$t \mapsto \mathbb{E}^{\mu^{*}}[F^{\pi} \mid \mathcal{F}_{t}]$ is a
martingale; the envelope is its essential supremum over policies, and its
evaluation at information states --- a Bellman recursion, not a
differential operator --- is the well-posed successor of the proposed
information-set operator.
\end{definition}

\begin{theorem}[finite well-posedness, comparison, exactness]
\label{dual-thm:envelope-finite}
In the finite setting --- disturbance set $\Omega$ finite, information
states $b$ drawn from a finite partition, admissible actions finite, data
rational --- define the one-step operator
\[
\Phi V(b) \;=\; \max_{a}\; \mathbb{E}^{\mu^{*}}\bigl[\, V(b') \,\big|\, b, a \,\bigr],
\]
where $b'$ is the posterior given $(b, a, y)$
the maximum over admissible actions of the prior-averaged next-belief
value. Then: (i) $\Phi$ is order-preserving,
$V \le W \Rightarrow \Phi V \le \Phi W$ pointwise; (ii)
$Q(0) = \Phi^{N} g$ is computed exactly by backward induction, in rational
arithmetic; and (iii) $Q(0) < 0$ certifies that for every admissible policy
the set of disturbance realizations defeating it carries positive
$\mu^{*}$-measure.
\end{theorem}
\begin{proof}
(i) If $V \le W$ pointwise, then for each $(b, a)$ the conditional
expectation of $V(b')$ is at most that of $W(b')$ (monotonicity of
expectation over the same kernel), and the maximum over the same action set
preserves the inequality. (ii) By induction on the remaining horizon:
$\Phi^{0} g = g$ is exact by definition; the induction step applies $\Phi$
once, and each evaluation is a finite weighted sum of previously computed
values, so rational data propagate exactly through the recursion.
Well-posedness --- existence and uniqueness of the recursion's output ---
is supplied by the same induction. (iii) $Q(0) < 0$ bounds every policy's
conditional expectation: for the optimal policy the expectation is
$Q(0) < 0$, and for every suboptimal $\pi$ the expectation is smaller
still, so Lemma~\ref{dual-lem:bridge} applies to each $\pi$ and yields the
positive-measure defeating set.
\end{proof}

\begin{proposition}[identification with the recourse certificate]
\label{dual-prop:recourse-id}
One step deep, with stored facet margins $\beta_{\ell}$, control kernels
$k_{\ell}(\sigma)$, prior $\lambda$ over the stored branches, and control
set $U$, the envelope evaluation of
Definition~\ref{dual-def:envelope} --- the policy's optimal ceiling on the
expected weighted facet sum --- is exactly the support-function certificate
\[
\Gamma_{h}(\lambda) \;=\; \sum_{\ell} \lambda_{\ell} \beta_{\ell}
\;+\; \int_{0}^{\tau} h_{U}\Big(\sum_{\ell} \lambda_{\ell} k_{\ell}\Big)
+ \sum_{\ell} \int_{\tau}^{T} h_{U}\big(\lambda_{\ell} k_{\ell}\big),
\]
with $h_{U}(g) = \max_{v \in U} g^{\top} v$ the support function of $U$,
of the obstruction calculus (its Section 3.7; complete proof in its
Supplementary Material, S1). The recourse certificate is the dynamic
envelope's one-step evaluation.
\end{proposition}
\begin{proof}
Variation of constants expands each realization of the weighted facet sum
into the stored $\beta$-term plus the window pairing (one common signal on
$[0, \tau)$) plus the branch pairings (separate post-reveal signals). The
supremum of a linear functional over the signal sets factorizes across the
time intervals, and each interval's supremum of a realized pairing
$g^{\top} v$ over $v \in U$ is the support value $h_{U}(g)$; collecting
terms gives exactly $\Gamma_{h}(\lambda)$ as the ceiling of every policy's
realized weighted facet sum. Negativity of the ceiling certifies by
Lemma~\ref{dual-lem:bridge}: the realized weighted facet sum of every policy is
at most $\Gamma_{h}(\lambda) < 0$, so some facet is violated on a
positive-measure set of realizations.
\end{proof}

\begin{corollary}[the three-branch instance, certified]\label{dual-cor:three-branch}
On the braking instance of the obstruction calculus (three branches,
weights $\lambda = (\tfrac{3}{8}, \tfrac{5}{16}, \tfrac{5}{16})$, stored
margin $\beta = \tfrac{16}{25} - \tau - 1$): the pooled window kernel
vanishes ($\sum_{\ell} \lambda_{\ell} n_{\ell} = 0$), each branch's
post-reveal pairing contributes its support value
$\lambda_{\ell} \int_{\tau}^{\tau+1} (\tau + 1 - \sigma)\, d\sigma =
\lambda_{\ell}/2$, and the envelope evaluation is
$\Gamma_{h}(\tau) = \tfrac{7}{50} - \tau$: negative exactly for
$\tau > \tfrac{7}{50} = 0.14$, where the bridge converts it into the
positive-measure defeating set for every blind-window policy;
hold-then-brake is safe exactly to $\tau = \tfrac{7}{50}$, so the
certificate is tight. The verification script evaluates the identity
exactly at $\tau \in \{\tfrac{1}{10}, \tfrac{7}{50}, \tfrac{1}{5}\}$.
\end{corollary}

\begin{conjecture}[the open residue in continuous time]\label{dual-conj:envelope}
The continuous-time residue of the envelope formulation: (i) an evolution
law for the adversarial prior along the review sequence --- a consistent
coupling of the one-step priors, in the spirit of martingale-transport
couplings (Beiglb\"ock, Henry-Labord\`ere, and Penkner, 2013) --- under which the envelope obeys a dynamic-programming equation
in continuous time; (ii) the limit of the information-state recursion under
partition refinement, which is where a comparison principle must live if a
differential statement is to be recovered at all; and (iii) value
attainment (the selection direction) in the general filtered setting. The
static identification of Proposition~\ref{dual-prop:recourse-id} is exact and is
not part of the conjecture.
\end{conjecture}

\section{The review-sequence envelope: exact couplings}

The envelope of the preceding section is one review deep. This section certifies
the two properties its continuous-time residue must preserve --- coupling
consistency and refinement monotonicity --- on exact review-sequence
instances; everything is rational arithmetic, regenerated by the script.

\subsection{The instance}

Two windows, one hidden branch. The branch is $s \in \{+1, -1\}$, prior
$\tfrac{1}{2}$ each; the state starts at $z_0 = 2$. At each review the
policy acts $u_1, u_2 \in \{-2, 0, 2\}$; between reviews the adversarial
disturbance acts $d_1, d_2 \in \{-1, 1\}$ (one consistent prior: the
product measure, $\tfrac{1}{2}$ per value per window). The state obeys
$z_1 = z_0 + s\, u_1 + d_1$ and $z_2 = z_1 + s\, u_2 + d_2$; the review
after window 1 observes $z_1$ exactly. The safety functional is the
expected two-sided margin
\[
F \;=\; \mathbb{E}_{s}\bigl[\, 2 - |z_{2}(s)| \,\bigr],
\qquad z_{2}(s) = z_0 + s(u_1 + u_2) + d_1 + d_2,
\]
bounded on the finite instance, so the bridge lemma applies verbatim.

\begin{proposition}[consistent coupling and the tower identity]
\label{dual-prop:tower}
The one-step adversarial priors couple to the product measure on
$(d_1, d_2)$, and for \emph{every} nonanticipative policy on the instance
--- all $3^{4} = 81$ policy trees over the reachable $z_1$ cells --- the
path-enumerated expectation and the iterated-conditional (tower)
evaluation of $\mathbb{E}[F]$ agree exactly. Each fixed policy's
conditional-expectation process along the review filtration is therefore
a martingale, and the envelope
$Q_0 = \operatorname*{ess\,sup}_{\pi} \mathbb{E}[F \mid \mathcal{F}_0]$
is computed by backward induction: $Q_0 = \Phi^{2} g$ equals the
policy-tree supremum ($= \tfrac{1}{2}$, at $u_1 = \pm 2$).
\end{proposition}
\begin{proof}
Consistency is by construction: the product measure has the one-step
priors as its marginals. For the tower identity, condition on
$\mathcal{F}_1$ (the review): within each $z_1$ cell the inner
conditional expectation is a finite weighted average, and the outer
average over the cells reconstructs the path sum because the cells
partition the paths --- an induction on the two-node filtration, exact
in rationals. The martingale property is the tower identity restated
along the filtration. Backward induction computes
$\Phi^{2} g$: $\Phi g$ at each cell is a maximum of three rational
expectations, and $\Phi$ is applied once more over $u_1$; the script
checks the result against the enumeration of all $81$ trees.
\end{proof}

\begin{proposition}[information design and recourse strictly pay]
\label{dual-prop:strict}
On the instance: the envelope value is $Q_0 = \tfrac{1}{2}$, attained by
$u_1 = -2$ (any $u_1 = \pm 2$ by symmetry) with the recourse
$u_2(z_1) = (0, -2, +2, +2)$ at $z_1 = (-1, 1, 3, 5)$; every open-loop
policy and the do-nothing policy attain only $0$. The strict gap
$\tfrac{1}{2} > 0$ is the certified content of the correction: the first
review's action \emph{designs the information} ($u_1 = -2$ makes $z_1$
separate the branch exactly), and the second \emph{exploits it}
(brake toward the origin on the revealed branch); at $z_1 = 5$ the
revealed state is unrecoverable and the cell value is $-1$.
\end{proposition}
\begin{proof}
The four singleton cells at $u_1 = -2$ carry values
$Q_1 = (1, 1, 1, -1)$ at weights $\tfrac{1}{4}$ each: $Q_0 = \tfrac12$.
The open-loop maximum is the better of nine fixed pairs, computed
exactly as $0$; the do-nothing pair $(0,0)$ also evaluates to $0$.
The cell values and the maximizing $u_2$ are direct rational
computations, reproduced by the script.
\end{proof}

\begin{proposition}[refinement is monotone]\label{dual-prop:refine}
On the three-state matching instance --- hidden $x \in \{1,2,3\}$, prior
$\tfrac13$ each, margin $-|x - a|$, action $a$ per observation block ---
the envelope is non-decreasing under partition refinement and strictly
increasing on the instance's final step:
$Q(\{\{1,2\},\{3\}\}) = Q(\{\{1\},\{2,3\}\}) = -\tfrac13
< 0 = Q(\{\{1\},\{2\},\{3\}\})$.
\end{proposition}
\begin{proof}
Every policy measurable under a coarser partition is measurable under a
refinement, so the supima are ordered. Strictness: under the fine
partition the policy matches each state exactly ($a = x$, value $0$);
under either coarse partition two states share an action and the best
shared action leaves one mismatch, value $-\tfrac13$. This monotone,
bounded family is the direction the partition-refinement limit of the
residue conjecture must take.
\end{proof}

The residue conjecture now reads sharper: the continuous-time coupling
law must reduce, on review sequences with finite disturbance sets, to
the product coupling certified here; and the information-state recursion
must converge along refinements with the monotonicity certified here.

\section{The general finite-sequence envelope}

The instance of the previous section is the $K = 2$ case of a general
statement. This section states the general theorem over finite review
sequences and proves it in full; the two-window instance is the $K = 2$ case, and the general argument uses only finite combinatorics and the linearity of expectation.

\subsection{Setup}

Fix a number of windows $K \ge 1$. Each window $k$ carries a finite
disturbance set $D_{k}$; the disturbance paths form
$\Omega = D_{1} \times \cdots \times D_{K}$, a finite set. A
\emph{consistent adversarial prior} is a product measure
$\mu^{*} = \mu^{*}_{1} \otimes \cdots \otimes \mu^{*}_{K}$ on $\Omega$;
consistency --- each marginal $\mu^{*}_{k}$ is the time-$k$ law of the one measure $\mu^{*}$ --- is exactly the requirement whose absence makes a backward recursion ill-posed. Reviews divide the sequence: the review
filtration is an increasing sequence of partitions
$\mathcal{P}_{1} \preceq \mathcal{P}_{2} \preceq \cdots \preceq
\mathcal{P}_{K}$ of $\Omega$, each $\mathcal{P}_{k}$ consisting of the
cells of the information available at review $k$ (so
$\mathcal{P}_{k}$ refines into $\mathcal{P}_{k+1}$), with
$\mathcal{P}_{K}$ the singleton partition; $\mathcal{P}_{1}$ may be the
trivial partition (a blind first window). Window $k$ has a finite action
set $A_{k}$; a \emph{nonanticipative policy} is
$\pi = (\pi_{1}, \dots, \pi_{K})$ with $\pi_{k}$ constant on the cells
of $\mathcal{P}_{k}$ --- the filtration carries the whole observation
history, so cell-constancy is adaptedness. The data are completed by a
bounded safety functional $G: \Omega \times A_{1} \times \cdots \times
A_{K} \to \mathbb{R}$ (positive values safe, per the sign convention of
this paper), rational throughout; the realized functional of $\pi$ is
$G^{\pi}(\omega) = G(\omega, \pi_{1}[\omega]_{1}, \dots,
\pi_{K}[\omega]_{K})$, where $[\omega]_{k}$ is the
$\mathcal{P}_{k}$-cell of $\omega$.

\begin{theorem}[general finite-sequence envelope]\label{dual-thm:general}
Under the setup above: \emph{(i) tower and martingale.} For every
nonanticipative policy $\pi$,
\[
\mathbb{E}^{\mu^{*}}[G^{\pi}] \;=\;
\mathbb{E}^{\mu^{*}}\bigl[\, \mathbb{E}^{\mu^{*}}\bigl[\, \cdots
\mathbb{E}^{\mu^{*}}[\, G^{\pi} \mid \mathcal{P}_{K} \bigr] \cdots
\mid \mathcal{P}_{2} \bigr] \mid \mathcal{P}_{1} \bigr],
\]
and $M^{\pi}_{k} = \mathbb{E}^{\mu^{*}}[G^{\pi} \mid
\mathcal{P}_{k}]$ is a $\mu^{*}$-martingale along the filtration.
\emph{(ii) Dynamic programming.} Define the node values $N$ on the
review tree --- nodes are pairs (cell, action prefix) --- by
$N_{K}\bigl(\{\omega\},\, a_{1:K-1}\bigr) = \max_{a \in A_{K}}
G(\omega, a_{1:K-1}, a)$ and, for $k < K$, the sum running over the
children $C' \subseteq C$ of $\mathcal{P}_{k+1}$ with the child weights
$w(C') = \mu^{*}(C')/\mu^{*}(C)$,
\[
N_{k}\bigl(C,\, a_{1:k-1}\bigr)
\;=\; \max_{a \in A_{k}} \;
\sum_{C'} w(C')\, N_{k+1}\bigl(C',\, a_{1:k-1}a\bigr),
\]
Then the envelope value $Q_{0} := \sup_{\pi}
\mathbb{E}^{\mu^{*}}[G^{\pi}]$ equals $N_{1}$ evaluated on the
$\mathcal{P}_{1}$-cells, averaged by their prior weights
$Q_{0} = \sum_{C \in \mathcal{P}_{1}} \mu^{*}(C)\,
N_{1}(C)$ (a single node if $\mathcal{P}_{1}$ is trivial); the supremum
is attained by a nonanticipative policy read off the maximizing
arguments of $N$. \emph{(iii) Refinement monotonicity.} If a second
filtration $\mathcal{P}'$ refines $\mathcal{P}$ stage-wise, then the
induced envelope values satisfy $Q_{0}(\mathcal{P}) \le
Q_{0}(\mathcal{P}')$. \emph{(iv) Exactness and the obstruction
direction.} All node values are rational; if $Q_{0} < 0$, then for
\emph{every} nonanticipative $\pi$ the set
$\{\omega: G^{\pi}(\omega) < 0\}$ carries positive $\mu^{*}$-measure.
\end{theorem}
\begin{proof}
\emph{(i)} Each $G^{\pi}$ is $\mathcal{P}_{K}$-measurable (a finite
function of the path and the cell-constant actions), and repeated
conditioning is an identity on finite probability spaces; the tower
formula is that identity read backwards, and the martingale property is
$M^{\pi}_{k} = \mathbb{E}[M^{\pi}_{k+1} \mid \mathcal{P}_{k}]$, which
is the defining property of conditional expectation applied twice.
\emph{(ii)} By backward induction on the stage. For $k = K$ the node
value is the best terminal action, and $G^{\pi}$ at a singleton cell is
$G(\omega, \cdot)$ evaluated at the realized prefix and $\pi_{K}$, so
$\max_{a_{K}} G = N_{K}$ pointwise. For the step, fix a
$\mathcal{P}_{k}$-cell $C$ and an action prefix, and split the policies
by their stage-$k$ action: every policy compatible with the prefix is a
choice $a \in A_{k}$ together with a continuation policy for the
conditioned sub-instance on $\{C\}$-successor cells, whose filtration
$\bigl(\mathcal{P}_{j}\bigr)_{j > k}$ restricted to the successors
refines to singletons at $K$. By (i) applied at the conditioning,
\[
\mathbb{E}^{\mu^{*}}[G^{\pi} \mid C] \;=\;
\sum_{C' \subseteq C} \frac{\mu^{*}(C')}{\mu^{*}(C)}\;
\mathbb{E}^{\mu^{*}}\bigl[G^{\pi} \mid C'\bigr],
\]
the conditional law of the successor cell given $C$ is the ratio of
prior masses (exogeneity: the prefix does not move $\mu^{*}$), and the
induction hypothesis --- applied to each conditioned sub-instance,
which has one fewer window --- identifies the supremum of the inner
conditional expectations over continuations of $(a_{1:k-1}, a)$ with
$N_{k+1}(C', a_{1:k-1}a)$. Taking the outer maximum over $a$ gives
$N_{k}$; summing over the $\mathcal{P}_{1}$-cells with their prior
weights and applying (i) once more gives $Q_{0} = \sum_{C}
\mu^{*}(C) N_{1}(C)$. Attainment: each $N$ node is a maximum over a
finite set; selecting a maximizing argument at every node and pasting
along the tree yields a nonanticipative policy --- cell-constancy holds
because the selection is made per node --- and its expectation is
exactly the recursion's value. \emph{(iii)} If $\mathcal{P}'$ refines
$\mathcal{P}$ stage-wise, every $\mathcal{P}$-cell-constant policy is
$\mathcal{P}'$-cell-constant, so the $\mathcal{P}'$-supremum ranges
over a superset of policies and cannot be smaller; the inequality for
the full sequences follows stage-wise from the one-step inclusion.
\emph{(iv)} The recursion's only operations are maxima and
$\mu^{*}$-weighted averages of rationals, so every $N$ is rational. If
$Q_{0} < 0$, then by (ii) every policy satisfies
$\mathbb{E}^{\mu^{*}}[G^{\pi}] \le Q_{0} < 0$, and the
expectation-to-realization bridge (Lemma~\ref{dual-lem:bridge}) converts strict negative expectation into a
positive-measure set of strictly negative realizations.
\end{proof}

\begin{remark}[the recourse certificate and the residue]\label{dual-rem:after-general}
Three consequences. \emph{First}, the identification with the recourse
certificate is the one-step case of Theorem~\ref{dual-thm:general}: with
stored facet margins the node recursion collapses to the
support-function evaluation $\Gamma_{h}$ (the identification
proposition), so the calculus's recourse
certificate is the envelope's $K$-node tree read at depth one.
\emph{Second}, the one-step priors $\mu^{*}_{k}$ do not determine the
coupling $\mu^{*}$: two consistent couplings with identical one-step
marginals can give different node values (the script certifies a
two-window example where the product coupling yields $0$ and the
comonotone coupling yields $+1$ on the same functional). The residue
conjecture's ``evolution law for the adversarial prior'' is therefore
precisely a \emph{coupling-selection} problem --- choose the consistent
coupling as a function of the state and review history --- which is the finite-space form of the martingale-transport structure recorded in the conjecture. \emph{Third}, refinement monotonicity (iii) is the
finite-space grounding of the partition-limit clause: the instance
family of the previous section, monotone and bounded, now carries a
general ordering, and the continuous-time residue is exactly the
question of which couplings and which refinement limits survive the
passage to continuous disturbance sets. The distributionally robust
interpolation of the probabilistic companion (Abaee, 2026, Probabilistic sufficiency) holds the
adversary to a mixture family: on the two-window instance its value is
exactly $\\rho$ at parameter $\\rho$ --- the interpolation contains the
coupling extremum at $\\rho = 1$, preserves the one-step marginals at
every $\\rho$, and leaves the coupling-selection difficulty orthogonal to
the mixture parameter (joint check archived).
\end{remark}

\section{Mechanized verification}

The duality results above are formalized in Lean~4 in \texttt{Formalizations/Minimax\_Dual} (\(21\) theorems and lemmas, importing only the
project prelude). The module builds as part of the project's \(60\)-job build, all of
which pass.

\noindent\textbf{No admitted gaps.} In Lean an unfinished proof appears as the axiom
\texttt{sorryAx} in the axiom footprint of every declaration depending on it. That
footprint was inspected: \texttt{sorryAx} is absent, and a source scan with comments
stripped finds no \texttt{sorry}, \texttt{admit} or \texttt{axiom} declaration anywhere in
the formalization layer. Declarations using classical reasoning depend only on \texttt{Classical.choice}, \texttt{propext} and \texttt{Quot.sound}, the standard axioms of Lean's classical logic.
\noindent\textbf{Verification provenance.} The project is pinned to Lean~4 in
\texttt{lean-toolchain} (currently \texttt{v4.34.1}). The mechanized checks reported in this
section were run under \texttt{v4.14.0}. On 2026-09-30 the source-level invariants were
re-verified against the current tree: a comment-stripped scan of all \texttt{.lean} files in
the project finds no \texttt{sorry}, \texttt{admit}, \texttt{axiom} or \texttt{constant}
declaration anywhere, and the declaration counts and named declarations quoted below were
re-counted and confirmed present. \textbf{The build has not been re-run since the toolchain
pin moved}, so the build-job figures below date from the \texttt{v4.14.0} run and should be
re-confirmed before submission.


\noindent\textbf{Scope.} What is machine-checked is the dual-certificate core stated over
an ordered-field interface, so it holds in any model of that interface. Numerical
instantiations, and any continuous-time statements, are not covered.

\section{Related work}
\label{dual-priorart}

The question this paper answers --- is there a single instrument that is safe against every
branch of an information set, and if not, what certifies that there is not --- has a long
history under other names. This section fixes the position against each of the four
neighbouring programmes. The short version: the \emph{question} is not new, and one of these
programmes already answers it in a different form; what is new is the \emph{form of the
answer} (a finitely supported measure with a tight support bound) and the \emph{precise
location of the boundary} at which that form stops existing.

\subsection{The honest baseline: this is the discriminating-kernel question}

The viability-theoretic treatment of two-player, state-constrained games is built on the
\emph{discriminating kernel} (Aubin, 1991; Aubin and Catt\'e, 2002; Cardaliaguet, Quincampoix
and Saint-Pierre, 1994, 2007): the largest closed subset of the constraint set that is a
discriminating domain, characterised as Victor's victory domain and, in the discrete and
continuous settings alike, computed by a descending-kernel algorithm. The discriminating
kernel is the two-player analogue of the viability kernel, and its algebraic characterisation
is as the largest, smallest, or \emph{unique minimax} --- bilateral --- fixed point of an
adequate map on pairs of subsets (Aubin and Catt\'e, 2002).

\textbf{This paper's obstruction is the emptiness of that object, and it should be read that
way.} The common safe-action set is nonempty exactly when the relevant discriminating domain
is nonempty, so Theorem~\ref{dual-thm:dual} is, in the vocabulary of that programme, a
certificate for the discriminating kernel being empty. That is a real collision with prior
art and is conceded here rather than argued around.

Four things are not in that programme, and they are what this paper claims.

\begin{enumerate}
\item \textbf{The certificate is a measure, not a set.} The discriminating-kernel programme
  returns a \emph{subset of state space}, characterised by a fixed-point property and computed
  by a set iteration. Theorem~\ref{dual-thm:dual} returns a \emph{probability measure} on the
  active boundary--disturbance bundle, finitely supported, whose existence is equivalent to
  emptiness. Membership in a kernel is thereby converted into the existence of a dual object
  that can be exhibited and checked without running the iteration.
\item \textbf{The support bound, and its tightness.} Proposition~\ref{dual-prop:sparse} gives a
  certificate supported on at most $k+1$ points, where $k$ is the \emph{control} dimension,
  and shows the bound tight. Note the scaling: the witness count is governed by the dimension
  of the instrument space, not of the state space. Nothing in the fixed-point characterisation
  of kernels yields a finite witness of this kind.
\item \textbf{The convexity boundary is located exactly where the kernel programme puts it,
  from the other side.} Cardaliaguet, Quincampoix and Saint-Pierre (1994, Proposition~2.3)
  show that under suitable convexity assumptions the discriminating kernel is
  \emph{convex}. Proposition~\ref{dual-prop:gap} here is the complementary negative: without
  convexity the duality fails outright, and an explicit two-action instance exhibits a strict
  minimax gap in which \emph{no} measure certifies. The two results are consistent and
  together say that the convexity hypothesis is load-bearing in the strongest possible sense
  --- it is simultaneously what makes kernels convex and what makes measure duality sound.
\item \textbf{The certificate is one-sided and observation-constrained, by design.} The
  discriminating-kernel programme computes a victory domain in both directions. This paper
  certifies only the negative direction (no common safe action exists) for an information set
  under partial observation. That restriction is what buys items~1 and~2; it is a narrowing,
  not an improvement, and is stated as such.
\end{enumerate}

\subsection{Farkas, and linear-programming duality}

The polyhedral case of Theorem~\ref{dual-thm:dual} recovers the $\ell_{1}$-normalised Farkas
multiplier pair exactly (Farkas, 1902), and the classical alternative is the mechanism on
which the whole certificate rests. Farkas' lemma is an alternative for a finite system of
linear inequalities: exactly one of a primal feasibility and a dual infeasibility certificate
holds.

The step beyond it is the measure-valued, non-polyhedral form. The classical alternative is
finite-dimensional and static; here the dual object ranges over probability measures on a
boundary--disturbance bundle, the emptiness being certified is that of a set of
\emph{simultaneously} safe actions across an information set rather than of a polyhedron, and
the recovery of the Farkas pair in the polyhedral case
(Proposition~\ref{dual-prop:recover}) is a consistency check on the generalisation, not the
generalisation itself.

\subsection{Isaacs, and the value-function programme}

The static picture is completed by the Isaacs drift condition (Isaacs, 1965). Where the
classical Hamilton--Jacobi--Isaacs theory produces a \emph{value} as a function of the state
under full observation, the object here is a one-directional refutation --- a proof that no
common instrument exists --- computed over an information set. The singleton case of
Theorem~\ref{dual-thm:dual} recovers the Isaacs drift with the adversarial \emph{distribution}
over disturbances load-bearing (Proposition~\ref{dual-prop:recover}); the distribution is what
carries the argument, and that is why the singleton case is not merely the classical
condition restated.

\subsection{Differential games under incomplete information, on the space of measures}

The closest single neighbour is the programme that lifts a zero-sum differential game with
\emph{imperfect} information about the state onto the space of probability measures.
Cardaliaguet and Quincampoix (2008) study the case in which the players know only a
distribution on the initial state, and establish existence of the value through a uniqueness
result for Hamilton--Jacobi--Isaacs equations stated on the space of measures. Both that
work and this one are measure-valued formulations of a game under incomplete information.

The aims differ. That programme establishes that a \emph{value exists}, by analytic means, in
an infinite-dimensional setting. This paper produces a \emph{finite certificate}: a
finitely supported measure on $k+1$ points, checkable in exact arithmetic, whose existence
refutes the common-action property. An existence-and-uniqueness result for a value does not
yield a finite witness, and the sparsity and tightness of
Proposition~\ref{dual-prop:sparse} have no analogue there. Conversely, nothing here establishes
the existence of a value, and the paper does not claim to.

\subsection{Minimax equality}

That minimax equality holds for quasi-concave--convex payoffs with appropriate compactness
and semicontinuity is classical (Sion, 1958), generalising von Neumann. The equality of
Theorem~\ref{dual-thm:dual} in the convex case can be read as an instance of that classical
circle of results, and this paper makes no claim to have originated the equality.

What the classical theorems do not supply is the sparse witness. Sion's theorem yields
\emph{existence}; it says nothing about the support size of a certifying measure, and the
$k+1$ bound of Proposition~\ref{dual-prop:sparse} --- together with the construction showing it is
saturated --- is not recoverable from it. Symmetrically, Proposition~\ref{dual-prop:gap} is
consistent with the classical theory: the minimax equality is precisely what
\emph{fails} once the convexity hypothesis is dropped.

\subsection{Verified numerics, and what is claimed}

The certificate is checked in exact rational arithmetic, with independent primal and dual
witnesses; the solver output is a confirmation, not a proof object. This is a statement about
the pipeline, not a claim over validated-integration methods, which address problems this
paper does not.

\medskip

\noindent \textbf{Claimed:} the measure form of the obstruction and its tight $k+1$ support
bound (Theorem~\ref{dual-thm:dual}, Proposition~\ref{dual-prop:sparse}); the two recoveries, Farkas and
Isaacs (Proposition~\ref{dual-prop:recover}); the exact location of the convexity boundary with an
explicit gap instance (Proposition~\ref{dual-prop:gap}); the dynamic, review-sequence and
finite-sequence envelope extensions with the tower identity and the strict-pay result
(Theorem~\ref{dual-thm:envelope-finite}, Proposition~\ref{dual-prop:tower},
Proposition~\ref{dual-prop:strict}, Theorem~\ref{dual-thm:general}); and exact machine verification
across all of it.

\noindent \textbf{Not claimed:} the discriminating-kernel question or its algorithm (Aubin;
Aubin and Catt\'e; Cardaliaguet, Quincampoix and Saint-Pierre), Farkas duality, the Isaacs
condition, minimax equality (Sion), or existence of a value under incomplete information
(Cardaliaguet and Quincampoix). The contribution is the \emph{certificate form} --- a
finitely supported, tightly bounded, exactly checkable measure witness --- and the
\emph{negative boundary result} that says precisely when that form ceases to exist.

\section{Conclusion}
\label{dual-conclusion}

The common-action obstruction says that no single instrument serves every branch of an
information set. Its polyhedral form carries a Farkas multiplier, and a Farkas multiplier is
already a certificate: a finitely supported witness that a system of linear inequalities is
infeasible. The question this paper asked is whether that is an accident of the polyhedral
case or the shadow of something general. It is the latter, and the general object has a
shape.

\textbf{What the obstruction is dual to.} For control-affine dynamics with a compact convex
control set, the common safe-action set is empty precisely when an adversarial probability
measure on the active boundary--disturbance bundle drives the expected inward drift strictly
negative against every control (Theorem~\ref{dual-thm:dual}). The obstruction is therefore not
merely a failure of a search to find a safe action; it is the existence of a \emph{witness} to
that failure, and the witness is a measure. This is the same movement that carries linear
programming from "the solver found nothing" to "here is a dual ray", and it carries the same
benefit: the certificate is checkable without re-running the search, and it localises the
failure on the active set.

\textbf{What governs the size of the witness.} The bound is the substantive point. The
certificate is supported on at most \(k + 1\) points, where \(k\) is the \emph{control}
dimension, and that bound is tight (Proposition~\ref{dual-prop:sparse}) --- it is attained, not
merely approached. Neither the state dimension, nor the cardinality of the disturbance set,
nor the resolution of any discretisation appears in it. So the complexity of certifying
nonviability is governed by how many independent controls the agent has, not by how large the
world it operates in is. A companion paper reaches an analogous conclusion for a different
pipeline: there the count is the information--time product rather than the input dimension
(Abaee, 2026, computational certification). The two are the same finding in different
languages, and the recurrence is the point --- \emph{obstruction certificates are small
objects whose size is set by the decision variables, not by the uncertainty.} That is what
makes them usable.

\textbf{What the two recoveries establish.} The polyhedral case recovers the
\(\ell_{1}\)-normalised Farkas multiplier and the game case recovers the Isaacs condition
(Proposition~\ref{dual-prop:recover}). Neither is a consistency check; each says that a classical
duality is a specialisation of this one, obtained by restricting the bundle. The measure form
is thus not a third thing alongside Farkas and Isaacs duality but the common generalisation
from which both drop out --- which is the sense in which it is the \emph{dual of the
obstruction} rather than a dual of one of its presentations.

\textbf{What the boundary result delimits.} The certificate form is not universal, and the
paper locates precisely where it stops rather than leaving that to be discovered: the convexity
boundary is identified exactly, with an explicit gap instance on the other side of it
(Proposition~\ref{dual-prop:gap}). Knowing when a certificate cannot exist is as much a part of the
contribution as knowing what it looks like where it does, because it converts a silent failure
into a diagnosed one.

\textbf{Scope, restated once.} Claimed here is the certificate form and the negative boundary,
together with the dynamic, review-sequence and finite-sequence envelope extensions and their
exact machine verification. Not claimed are the discriminating-kernel question and its
algorithm, Farkas duality, the Isaacs condition, minimax equality, or the existence of a value
under incomplete information --- all of which are assumed, cited, or deliberately outside the
frame. The contribution is a \emph{form}: a finitely supported, tightly bounded, exactly
checkable measure witness to nonviability, and a statement of when that form ceases to exist.

\section{Cross-part conclusion}
\label{conclusion}

The two parts above build obstruction certificates by two different routes --- a finite linear
program over moment relaxations, and a measure-dual characterisation of emptiness --- and they
arrive at the same structural conclusion. That recurrence is the paper's central result.

\emph{The shared principle.} \textbf{Obstruction certificates are small objects whose size is set
by the decision variables, not by the uncertainty.} Part I: whenever a positive certificate exists,
one can be chosen with at most $r+1$ safety rows, $r$ the information--time rank. Part II: the
witness measure is supported on at most $k+1$ points, $k$ the control dimension, and the bound is
\emph{tight} --- attained, not merely approached. In neither bound does the state dimension, the
cardinality of the disturbance set, or the resolution of any discretisation appear. The two counts
are not the same number and are not claimed to be: they count different decision variables in
different pipelines. What is shared is the \emph{form} of the dependence. Certifying that something
cannot be done costs a number of degrees of freedom set by what the agent may decide, and is
otherwise insensitive to how much the agent does not know.

\emph{Why the bound is substantive rather than decorative.} A bound is only interesting if it can
fail, and Part I exhibits the failure of the natural alternative. A Helly-type statement --- at most
$m+1$ compatible states suffice to witness nonviability --- is valid for convex common-action sets
in $\mathbb{R}^m$ and is \emph{false} for a common blind control function: for every $q$ there is a
system with scalar input and $q+1$ indistinguishable modes in which every proper subbelief is
viable while the full belief is not, so every certificate must involve all $q+1$ modes. The
correct dimension therefore counts independent \emph{temporal and informational} control decisions,
not inputs. This is not a technical footnote: it decides whether certificate-size claims for
review intervals, delayed sensing, and adaptive observation are true, and it is the reason the
information--time product --- rather than any input dimension --- is the quantity that appears in
Part I's bound.

\emph{What the obstruction is, and why that matters.} Part II's contribution is ontological rather
than algorithmic: the obstruction is not \emph{merely a failure of a search to find a safe action};
it is the existence of a \emph{witness} to that failure, and the witness is a measure. This is the
same movement that carries linear programming from ``the solver found nothing'' to ``here is a dual
ray,'' and it carries the same benefit: the certificate is checkable without re-running the search,
and it localises the failure on the active set. Read together with Part I, the two parts therefore
give the negative verdict the same standing that a positive one enjoys --- a nonviability claim
becomes an artefact a reader can verify, not an assertion a reader must trust.

\emph{Where the equivalence fails, stated plainly.} The measure characterisation requires convexity.
A two-action instance exhibits a strict minimax gap: without convexity the equivalence fails and
\emph{no measure certifies}. Part I's guarantee is correspondingly one-directional in general ---
nonviability is eventually detected along admissible refinement sequences, but a failure to detect
at finite refinement is not a proof of viability. Both parts state their own scope delimitations
rather than letting the shared principle imply more than it does.

\emph{The two recoveries, and what they license.} The polyhedral case recovers the
$\ell_1$-normalised Farkas multiplier exactly, and the singleton case recovers the Isaacs drift
with the adversarial distribution over disturbances load-bearing. So the classical objects --- Farkas
infeasibility certificates and Isaacs minimax drift conditions --- are not rival frameworks but
specialisations of one duality, and a practitioner already using either is already using this one.
What the unification licenses is the transfer of the checkability property: a certificate that can
be verified without re-running the search, whose size is known in advance, and which localises the
failure it reports.

\subsection*{References}
\label{references}
Abaee, A., 2026. An obstruction calculus for viability under incomplete
observation. Manuscript submitted for publication.

Abaee, A., 2026. Exact audits of worked systems for the obstruction
calculus: counts, witnesses, and monitoring design rules. Manuscript
submitted for publication.

Abaee, A., 2026. Exact audits of worked systems for the obstruction calculus:
counts, witnesses, and monitoring design rules (submitted).

Abaee, A., 2026. Exact belief-state computation at scale II: the
four-parameter cube and the antichain census. Manuscript submitted for
publication.

Abaee, A., 2026. Minimax dual certificates for the common-action
obstruction. Manuscript submitted for publication.

Abaee, A., 2026. Probabilistic sufficiency for the obstruction calculus:
belief-state safety values under partial observation (submitted).

Abaee, A., 2026. Robust viability of the 2J3KL limit reference point
under a surplus-production map: policy scoring, expansion, and when catch
cannot help. Zenodo. \url{https://doi.org/10.5281/zenodo.22552060}

Anderson, E.J., Nash, P., 1987. Linear Programming in Infinite-Dimensional
Spaces: Theory and Applications.

%% ------------------------------------------------------------------
%% Added 2026-09-29 to support section \ref{priorart}. All VERIFIED against
%% publisher records (journal, volume, pages, year confirmed):
%%   Saint-Pierre 1994, Appl. Math. Optim. 29, 187-209 (doi 10.1007/BF01204182)
%%   Maidens, Kaynama, Mitchell, Oishi & Dumont 2013, Automatica 49(7), 2017-2029
%%   Mitchell, Bayen & Tomlin 2005, IEEE Trans. Automat. Control 50(7), 947-957
%%   Lasserre 2001, SIAM J. Optim. 11(3), 796-817
%%   Parrilo 2003 - SOS relaxations (survey/volume; page numbers not confirmed)
%%   Calafiore & Campi 2005, Math. Program. 102(1), 25-46
%%   Campi & Garatti 2008, SIAM J. Optim. 19(3), 1211-1230
%%   Campi & Garatti 2011, J. Optim. Theory Appl. 148(2), 257-280
%%   Campi, Garatti & Prandini 2009, Annu. Rev. Control 33(2), 149-157
%% NOTE: the Automatica paper is Maidens et al. (Mitchell is THIRD author, not
%% first). Cited correctly. Do not "correct" it to Mitchell et al.
%% ------------------------------------------------------------------
Calafiore, G. and Campi, M.C., 2005. Uncertain convex programs: randomized solutions and
confidence levels. \emph{Mathematical Programming} \textbf{102}(1), 25--46.

%% ------------------------------------------------------------------
%% Added 2026-09-29 to support section \ref{priorart}. All VERIFIED against
%% publisher records (journal, volume, pages, year confirmed):
%%   Aubin 1991, Viability Theory, Birkhauser Boston
%%   Aubin & Catte 2002, Set-Valued Anal. 10, 379-416 (doi 10.1023/A:1020667819804)
%%   Cardaliaguet, Quincampoix & Saint-Pierre 1994, RAIRO Modelisation Math.
%%     Anal. Numer. (M2AN) 28(4), 441-461  -- the discriminating-kernel ALGORITHM
%%     paper; its Prop 2.3 is the convexity-of-kernel result this paper's
%%     prop:gap is the complement of.
%%   Cardaliaguet, Quincampoix & Saint-Pierre 2007, in Advances in Dynamic Game
%%     Theory, Ann.

Aubin, J.-P. and Catt\'e, F., 2002. Bilateral fixed-points and algebraic properties of
viability kernels and capture basins of sets. \emph{Set-Valued Analysis} \textbf{10},
379--416.

Aubin, J.-P., 2001. Viability kernels and capture basins of sets under differential
inclusions. \emph{SIAM Journal on Control and Optimization} \textbf{40}(3), 853--881.

Aubin, J.-P., 1991. Viability Theory.

Beiglb\"ock, M., Henry-Labord\`ere, P., Penkner, F., 2013.
Model-independent bounds for option prices --- a mass transport approach.

Birkh\"auser, Boston.

Birkh\"auser Boston, Boston, MA.

Calafiore, G. and Campi, M.C., 2006. The scenario approach to robust control design.
\emph{IEEE Transactions on Automatic Control} \textbf{51}(5), 742--753.

Campi, M.C. and Garatti, S., 2011. A sampling-and-discarding approach to chance-constrained
optimization: feasibility and optimality. \emph{Journal of Optimization Theory and
Applications} \textbf{148}(2), 257--280.

Campi, M.C. and Garatti, S., 2008. The exact feasibility of randomized solutions of uncertain
convex programs. \emph{SIAM Journal on Optimization} \textbf{19}(3), 1211--1230.

Campi, M.C., Garatti, S. and Prandini, M., 2009. The scenario approach for systems and
control design. \emph{Annual Reviews in Control} \textbf{33}(2), 149--157.

Cardaliaguet, P. and Quincampoix, M., 2008. Deterministic differential games under
probability knowledge of initial condition. \emph{International Game Theory Review}
\textbf{10}(1), 1--16.

Cardaliaguet, P., Quincampoix, M., Saint-Pierre, P., 1999. Set-valued
numerical methods for optimal control and differential games, in:
Stochastic and Differential Games: Theory and Numerical Methods,
pp.~177--247.

Cardaliaguet, P., Quincampoix, M. and Saint-Pierre, P., 1994. Some algorithms for differential
games with two players and one target. \emph{RAIRO -- Mod\'elisation Math\'ematique et Analyse
Num\'erique} \textbf{28}(4), 441--461.

Cardaliaguet, P., Quincampoix, M. and Saint-Pierre, P., 2007. Differential games through
viability theory: old and recent results. In: J\o rgensen, M., Quincampoix, M. and Vincent,
T.L. (eds), \emph{Advances in Dynamic Game Theory}, Annals of the International Society of
Dynamic Games vol.\ 9, pp.\ 3--35. Birkh\"auser Boston.

Chv\'{a}tal, V., 1983. Linear Programming. A Series of Books in the
Mathematical Sciences. W.~H.

Farkas, J., 1902. Theorie der einfachen Ungleichungen. J.\ Reine Angew.
Math. 124, 1--27.

Finance and Stochastics 17, 477--501.

Freeman, New York.

Helly, E., 1923. \"Uber Mengen konvexer K\"orper mit gemeinschaftlichen
Punkten. J.~Deutsche Math.-Ver. 32, 175--176.

Hettich, R., Kortanek, K.O., 1993. Semi-infinite programming: theory,
methods, and applications. SIAM Rev. 35(3), 380--429.

Huangfu, Q., Hall, J.A.J., 2018. Parallelizing the dual revised simplex
method. Math. Program. Comput. 10(1), 119--142.

Isaacs, R., 1965. Differential Games: A Mathematical Theory with
Applications to Warfare and Pursuit, Control and Optimization. Wiley,
New York.

ISDG 9, 3-35, Birkhauser Boston
%%   Cardaliaguet & Quincampoix 2008, Int. Game Theory Rev. 10(1), 1-16 --
%%     value via HJI ON THE SPACE OF MEASURES under probability knowledge of the
%%     initial condition. The single closest neighbour.
%%   Sion 1958, Pacific J. Math. 8(1), 171-176 (doi 10.2140/PJM.1958.8.171)
%% NOTE ON AUTHOR NAMES: it is Aubin and CATTE (accented), not "Catte".
%%   Do not strip the accent.
%% NOTE ON A DROPPED CLAIM: an earlier roadmap note attributed to Aubin (2001)
%%   the "complement of the kernel -- Poincare's shadow". That attribution could
%%   NOT be verified and is NOT asserted here. Aubin (2001) is cited only for
%%   the verified fixed-point characterisation of kernels and capture basins.
%% ------------------------------------------------------------------
Aubin, J.-P., 1991. \emph{Viability Theory}. Systems and Control: Foundations and
Applications.

Lasserre, J.B., 2001. Global optimization with polynomials and the problem of moments.
\emph{SIAM Journal on Optimization} \textbf{11}(3), 796--817.

Maidens, J.N., Kaynama, S., Mitchell, I.M., Oishi, M.M.K. and Dumont, G.A., 2013. Lagrangian
methods for approximating the viability kernel in high-dimensional systems. \emph{Automatica}
\textbf{49}(7), 2017--2029.

Milanese, M., Norton, J., Piet-Lahanier, H., Walter, E. (Eds.), 1996.
Bounding Approaches to System Identification.

Mitchell, I.M., Bayen, A.M. and Tomlin, C.J., 2005. A time-dependent Hamilton--Jacobi
formulation of reachable sets for continuous dynamic games. \emph{IEEE Transactions on
Automatic Control} \textbf{50}(7), 947--957.

Parrilo, P.A., 2003. Semidefinite programming relaxations for semialgebraic problems.
\emph{Mathematical Programming} \textbf{96}(2), 293--320.

Plenum Press, New York.

Saint-Pierre, P., 1994. Approximation of the viability kernel. \emph{Applied Mathematics and
Optimization} \textbf{29}, 187--209.

Sion, M., 1958. On general minimax theorems. Pacific J. Math. 8(1),
171--176.

Sion, M., 1958. On general minimax theorems. \emph{Pacific Journal of Mathematics}
\textbf{8}(1), 171--176.

Sion, M., 1958. On general minimax theorems. Pacific J.\ Math. 8, 171--176.

Veliov, V.M., 1993. Sufficient conditions for viability under imperfect
measurement. Set-Valued Anal. 1, 305--317.

Wiley, Chichester.

\subsection*{Declarations}

\subsection*{Funding}
None.

\subsection*{Competing interests}
None.

\subsection*{Data availability}
No external data were used.

\subsection*{Code availability}
The verification scripts
\url{paper2_computational_certification_v20_verification.py},
\url{paper2_nonstandard_verification.py},
\url{paper2c_lp_campaign.py}, and the \texttt{viacert} package are publicly
available at \url{https://github.com/MIKEAA2020/general-sustainability}
(folder \texttt{arena agent 1/paper rewrites/latex}, at the commit
pinned to this article's submission); they reproduce all exact values,
witnesses, bounds, and the solver campaign verbatim.

\subsection*{AI declaration}
GLM (Z.ai), Qwen (Alibaba Cloud) and DeepSeek AI assisted with drafting
and iterative review. The author reviewed and edited outputs and takes
responsibility for the final work.
\section*{Declarations}
\textbf{Funding.} None.\quad \textbf{Competing interests.} None.\quad
\textbf{Data availability.} No external data were used.\quad
\textbf{Code availability.} The verification script
\texttt{minimax\_dual\_certificates\_v6\_verify.py} is publicly available
at \url{https://github.com/MIKEAA2020/general-sustainability} (folder
\texttt{arena agent 1/paper rewrites/latex}); it reproduces every check family verbatim. The cross-paper certificate
exchange parser (\texttt{certificate\_exchange\_v1.py}) validates the
archived dual certificate of this paper and of the certification companion
exactly, in either direction.\quad
\textbf{AI declaration.} GLM (Z.ai), Qwen (Alibaba Cloud) and DeepSeek AI
assisted with drafting and iterative review. The author reviewed and
edited all outputs and takes responsibility for the final work.

\end{document}
